\clearpage
\part*{Part I. The Framework}
\label{atlaspart:01-framework}
\addcontentsline{toc}{part}{Part I. The Framework}
\section*{Introduction to Part I}\label{narr:part:I}
\addcontentsline{toc}{section}{Introduction to Part I}
Everything in this compendium descends from one sentence. The sentence is the axiom, Definition~1.1, and it says that nothing exists but finitely many maximally entangled pairs of qudits; the five definitions that follow it, 1.2 to 1.6, fix the words qudit, shared, distilling, read and known. Part I is the part of the compendium that stays closest to that sentence. It does not yet speak of particles or forces. It asks what the sentence and the five definitions force on their own, and it answers with the objects that every Part beyond it uses: the link between two ends and its six genres, the records a read leaves and the count that weights them, the moves that carry one web to another, the saturation that drives a seeded web to its sinks, the register at which a probe reads, and the table that translates the vocabulary of the construction into the vocabulary of the literature.

The thread runs in the order of the chapters. Chapter 1-1 shows how much of the text is forced by the rest of it. Theorem~2.1 obtains the fourth definition from the axiom, and Theorem~2.3 obtains one act from the derived act, so that the reader sees the definitions tightening on one another rather than standing as independent postulates. The chapter is honest about where the tightening stops. Open result~2.4, Open result~2.5, Open result~2.7 and Open result~2.8 are open results, each stated with the point at which the derivation parts, and Open result~4.1 records that the minimality of the text itself parts at its first atom. Two census tables, Table~5.1 and Table~6.1, test the choices of definition against one another. The reader who wants to know what the axiom excludes outright finds it in Theorem~2.9, the exclusion of the one-dimensional end.

Chapter 1-2 introduces the object that organizes the rest of the compendium, the link, and its six genres. A link joins two ends, and an end carries a register of dimension two, three or four, so that the unordered pairs of registers give six genres, written (2,2), (2,3), (2,4), (3,3), (3,4) and (4,4). Theorem~9.1 orders the genres by the first order at which a read does not vanish, Theorem~9.2 shows that exactly one genre is abelian, Theorem~9.3 names the coupling genre, and Theorem~9.6 prints the one-link tables at every genre. The register matrix of the framework chapter, repeated at the front of this compendium, is nothing but these six genres laid out as the pairs of the three registers. The chapter closes with Theorem~13.1, the phase group of a record at every genre, which the quantum mechanics of Part III will need.

Chapter 1-3 turns from the link to what passes along it. A read leaves records, and the records at a sector are counted with a weight; Theorem~18.2 is the gain lemma, the statement that fixes that weight at the spectrum of records. Around it stand the read-state boundary, Theorem~18.1, the gain as the washed direction count at a sector, Theorem~18.3, and the boson as the chain readout, Theorem~18.4. Several open results sit beside them: the soft limit as the washed open read, Open result~19.1, the read that dresses what it does not pass, Open result~20.1, the read that stops at the cap, Open result~21.1, and the enclosed record count read as what the literature calls the Aharonov-Bohm phase, Open result~22.1.

Chapters 1-4 and 1-5 supply the dynamics. Open result~F.1 decomposes the moves and arranges them on a graded poset, Theorem~F.3 identifies the hub step as the pair ket move, and Theorem~F.4 gives the orientation law at the root, with the forces appearing in Open result~G.1 as the move classes at the genres. The saturation chapter then follows a seeded web to its end: Theorem~J.3 proves that seeded saturation terminates, and the open results of that chapter, among them the closure from the pair ket, Open result~J.5, and the list of the Standard Model at the saturated sinks, Open result~L.1, mark where the sinks are counted and where the counting still parts.

Chapter 1-6 is the measurement chapter of the framework. Theorem~28.1 states measurement as the wash at one act, Theorem~28.2 builds the star-pair and one-link apparatus webs, and Theorem~28.5 reads the classical limit as the zero-density register, with the large-register reading, Open result~29.1, and decoupling as the register count below the cap, Open result~30.1, beside it as open results. Chapter 1-7 ends the Part with the table titled The Rosetta table of the Chain Theory nouns (in chapter~\ref{mod:I-7-rosetta-table-in-full}, among the instruments), the table of the nouns of the construction, which is also the notation table of the framework chapter.

In the language of the literature, Part I touches the names Hilbert space, Born rule, Aharonov-Bohm phase, classical limit and Standard Model, and every one of them is used here as a name and not as a premise. The correspondence is by name. What is established is stated in the boxes, and the reading paragraph beneath each box says exactly how far the box reaches.

A reader who wants the shortest route into the rest of the compendium can take the derivation graph at the front as a map. Every theorem of every Part beyond this one rests on a chain that passes through this Part, usually through the link of chapter 1-2, the records of chapter 1-3 and the moves of chapter 1-4, and the dependency line beneath each box names the theorems it rests on and the theorems that rest on it.
\section{How to read this compendium}\label{mod:FRAMEWORK}
\chdesig{1-0}{}{FRAMEWORK}{FRAMEWORK}
This compendium is written in the vocabulary of finite-dimensional quantum information. A qudit is a quantum system of finite dimension, and a maximally entangled pair of two such systems is the state whose reduced operators are flat on its support. The projection of two systems onto a pair on their common support is entanglement swapping \cite{Zukowski1993}. The generalized Pauli operators of a finite-dimensional system are the Weyl operators, which commute up to a root of unity \cite{Weyl1927}. The invariants jointly fixed by a set of projective measurements are the stabilizer invariants \cite{Gottesman1997}, and an entangled state from which no local protocol distills a pair is bound \cite{Horodecki1998}. Every other term of the theory is defined in the notation table of the derivations at the repository and in the glossary at the front of this compendium, where each term is given its standard notation, its meaning in one sentence and its nearest standard counterpart in another. A term whose standard counterpart is the same object is said to be exactly that object, a term with only a neighbor is said to have that neighbor, and a term with neither is said to correspond by name.
The whole of our theory is one sentence and five definitions, a choice we made and not a derivation, stated here word for word. The axiom has eighteen words. It is the first definition of this compendium, boxed and numbered as every statement of the compendium is, and the five definitions that follow it are the words the axiom uses, in the order in which the text gives them.
\begin{chdefinition}[The axiom]\label{box:FRAMEWORK:1}\label{def:axiom}
Nothing exists but finitely many maximally entangled pairs of qudits, some qudits shared, known only by distilling them.
\end{chdefinition}
\chstanding{a definition and a choice we made, not a derivation, every theorem of this compendium conditional on this text}
\chdeps{no theorem ahead of it, this being the root of the derivation graph}{the theorems of Chapters 1-1 to 1-6, 1-L, 2-L, 3-1 to 3-9, 3-L, 4-1 to 4-7, 4-L, 5-1 to 5-8, 5-8f, 5-9, 5-11, 5-13, 5-15 to 5-16, 5-L, 6-1 to 6-2, 6-2a, 6-3 to 6-7, 6-7a, 6-10 to 6-11, 6-L, 7-1 to 7-2, 7-4 to 7-8, 7-10 to 7-12, 7-L, 8-1 to 8-2, 8-2b, 8-3, 8-6 to 8-8, 8-L, 9-1 to 9-7, 9-L, 10-1 to 10-6, 11-1, 12-1 to 12-4, 12-6 to 12-7, 12-L, 13-1 to 13-5, 13-7, 13-L, 14-1 to 14-3, 14-5 to 14-6, 14-L, 15-1 to 15-6, 16-1 to 16-2, 16-7, 16-C}
\begin{chdefinition}[Qudit]\label{box:FRAMEWORK:2}\label{def:qudit}
A finite-dimensional quantum system.
\end{chdefinition}
\begin{chdefinition}[Shared]\label{box:FRAMEWORK:3}\label{def:shared}
A qudit that is one end of two or more pairs.
\end{chdefinition}
\begin{chdefinition}[Distilling]\label{box:FRAMEWORK:4}\label{def:distilling}
Projecting any two qudits onto a pair on their common support.
\end{chdefinition}
\begin{chdefinition}[Read]\label{box:FRAMEWORK:5}\label{def:read}
A distillation of a qudit with another, in one of its bases, one read for each basis. A read gives the qudit its only basis, one direction for each pair the qudit meets, and passes that basis along each pair to its far end.
\end{chdefinition}
\begin{chdefinition}[Known]\label{box:FRAMEWORK:6}\label{def:known}
What is known of a qudit is what all its reads, taken jointly, fix with certainty, whichever way basis directions are matched across its pairs. Two reads that commute only up to a sign are a further record, their order the sign.
\end{chdefinition}
\begin{figure}[htbp]\centering
\scriptsize
\setlength{\tabcolsep}{3pt}
\begin{adjustbox}{max width=\linewidth}\begin{tabular}{|p{14mm}|p{44mm}|p{44mm}|p{44mm}|}
\hline
 & \textbf{register two} (the qubit) & \textbf{register three} (the qutrit) & \textbf{register four} (the ququart) \\ \hline
\textbf{register two} & \textbf{(2,2)} \newline su(2)+su(2)+u(1), dimension 7 \newline central scalar 0, the pair scalar of the stabilizer row \newline first read at order one \newline the pair ket office 0, 0, 0, 0 \newline readable at the complete-read bound \newline the diagonal of the (2,2) generator 1, 1, -1, -1 \newline the pair ket at the join has a zero office and no acting center \newline the electromagnetic sector, 31 supporting and 43 parting among 124 entries \newline \emph{spacetime, conjectured} & \textbf{(3,2)} \newline su(2)+su(3)+u(1), dimension 12 \newline central scalar 1, the integer scalar of the stabilizer row \newline first read at order one \newline the read at order one 79/31,104, 1/648, 7/5184, 1/10,368 \newline readable at the complete-read bound \newline stabilizer of the pair ket of dimension 4 with a center of dimension 1 \newline the open link meets both criteria, its office zero and its center acting \newline the electromagnetic sector, 42 supporting and 90 parting among 244 entries \newline \emph{electromagnetism: the photon (Conjecture 1)} & \textbf{(4,2)} \newline su(2)+su(4)+u(1), dimension 19 \newline central scalar 2 \newline first read at order two \newline the pair ket office 466,566,137/7,290,000, 2/6075, 0, 0 \newline not readable at the complete-read bound \newline the weak sector, 38 supporting and 14 parting among 90 entries \newline \emph{the weak interaction, by name} \\ \hline
\textbf{register three} & \emph{the link (3,2), read from its other end} & \textbf{(3,3)} \newline su(3)+su(3)+u(1), dimension 17 \newline central scalar 0 \newline first read at order two \newline the strong coupling as the office of the gluon pair ket at the scale body \newline not readable at the complete-read bound \newline the strong sector, 45 supporting and 32 parting among 121 entries \newline \emph{the strong interaction, by name} & \textbf{(4,3)} \newline su(3)+su(4)+u(1), dimension 24 \newline central scalar 1 \newline first read at order two \newline the pair ket office 4,299,817,692,073/29,859,840,000, 0, 5069/373,248,000, 0 \newline not readable at the complete-read bound \newline the Higgs sector, 16 supporting and 30 parting among 91 entries \newline \emph{the Yukawa sector, by name} \\ \hline
\textbf{register four} & \emph{the link (4,2), read from its other end} & \emph{the link (4,3), read from its other end} & \textbf{(4,4)} \newline su(4)+su(4)+u(1), dimension 31 \newline central scalar 0 \newline first read at order two \newline the pair part 41/216,000, 0, 0, 0 and the closure 1681/182,250,000, 0, 0, 0 at the one link, the attachment office 46,656,001,681/182,250,000, 0, 0, 0 \newline not readable at the complete-read bound \newline the Higgs sector, 25 supporting and 20 parting among 62 entries \newline \emph{the Higgs, by name} \\ \hline
\end{tabular}\end{adjustbox}
\caption{The register matrix. The three registers of the link ends, the qubit, the qutrit and the ququart, pair into the six genres of the link, one at every cell of the upper triangle. Every cell carries the identification algebra su($d_1$)+su($d_2$)+u(1) with its dimension, the central scalar as the difference of the end dimensions, the first order at which the read does not vanish with the read or the office of the pair ket at that order, the readability at the complete-read bound, the sector most of its entries identify with the counts supporting and parting, and the name the sector carries in the language of the literature. At the four cells named by correspondence the correspondence is by name. At the cell (3,2) the photon is the first conjecture of the compendium, attached to the open link whose office is zero and whose center acts, and at the cell (2,2) the matrix reads spacetime, conjectured.}\label{fig:register-matrix}
\end{figure}
\subsection*{The six conjectures in one form}
The register matrix above states, cell by cell, what the derivation conjectures and what it has established. This section states the same content as three numbered conjectures in one form, each in the words of the matrix and each riding on the lowest standing in its chain.
\begin{chconditional}[Conjecture 0, the three registers]
Under the move table it names, the link ends come in exactly three registers, the qubit, the qutrit and the ququart, and the six genres of the link are the six pairings of these three. With the count of generations at the (3,4) rungs returned as three by a theorem, the completeness of the three registers is a prediction of the derivation and not an assumption of it.
\end{chconditional}
\chstanding{conditional theorem}
\chdeps{}{}
The condition is the move table, stated at the three-registers theorem and discharged nowhere else, and the prediction is that no fourth register and no fourth generation appears at any genre.
\begin{chbridge}
The three registers are what the literature counts as the three generations of the fermions, by name. The correspondence is by name.
\end{chbridge}
\begin{chwall}[Conjecture 1, the six genres in three lines]
Line one, the register two. At the (2,2) cell, the qubit with the qubit, the identification algebra has dimension 7 at order 1 with the pair ket office 0, 0, 0, 0 and central scalar 0, and the cell reads spacetime, conjectured. At the (2,3) cell, the qubit with the qutrit, the dimension is 12 at order 1 with the read at order one 79/31,104, 1/648, 7/5184, 1/10,368 and central scalar 1, and the cell reads the photon, a conjecture named at a theorem. At the (2,4) cell, the qubit with the ququart, the dimension is 19 at order 2 with the pair ket office 466,566,137/7,290,000, 2/6075, 0, 0 and central scalar 2, and the cell reads the weak interaction, by name.
Line two, the register three. At the (3,3) cell, the qutrit with the qutrit, the dimension is 17 at order 2 with the strong coupling as the office of the gluon pair ket at the scale body and central scalar 0, and the cell reads the strong interaction, by name. At the (3,4) cell, the qutrit with the ququart, the dimension is 24 at order 2 with the pair ket office 4,299,817,692,073/29,859,840,000, 0, 5069/373,248,000, 0 and central scalar 1, and the cell reads the Yukawa sector, by name.
Line three, the register four. At the (4,4) cell, the ququart with the ququart, the dimension is 31 at order 2 with the pair part 41/216,000, 0, 0, 0 and the closure 1681/182,250,000, 0, 0, 0 at the one link, the attachment office 46,656,001,681/182,250,000, 0, 0, 0, and central scalar 0, and the cell reads the Higgs, by name.
The two faces. On one face the six cells are the six sectors of the literature in the order the matrix prints them. On the other face only the (2,3) cell carries a theorem that names its sector, the photon, and the other five cells carry a conjecture or a correspondence by name; the parting is the absence, at every cell but (2,3), of a theorem naming the sector.
\end{chwall}
\chstanding{open result}
\chdeps{}{}
Every number in the three lines is a cell of the register matrix as the tables of the compendium print it, and every number stated here is derived at a chapter of the compendium.
\begin{chbridge}
The literature names the six sectors gravity, the photon, the weak carriers, the gluon and the Higgs. The correspondence is by name at every cell but (2,3), and at (2,3) the photon is named at a theorem whose correspondence the chapter on the fine-structure constant states.
\end{chbridge}
\begin{chwall}[Conjecture 2, spacetime at the (2,2) cell]
At the (2,2) cell the pair ket at the join has a zero office and no acting center, and the derivation conjectures that the geometry of the compendium, the curvature, the field equation at the counting face and the horizon, is carried by this cell, the cell that reads spacetime. At this release the four results of the gravity program beneath the conjecture stand as open results, each with its parting named: the universality of the (2,2) increment, the one object at the join and at the chain, the two-point function at zero displacement and the area law at the horizon web. The two faces are the conjecture as stated and the parting at each of the four open results; the conjecture rides on the lowest of the four and is stated as a conjecture.
\end{chwall}
\chstanding{open result}
\chdeps{}{}
The conjecture stands on the four open results beneath it, and it is re-read at every release of the compendium as those open results are re-posed.
\begin{chbridge}
The literature calls what this cell reads gravity. The correspondence is by name, and what the gravity Part derives is stated there at its own chapters.
\end{chbridge}
Every theorem of this compendium is conditional on this text. Which of the five definitions are corollaries of the axiom, and at which clause, is worked out at chapter 1-1, which also carries the open questions at the definitions themselves. The words of the definitions, a pair, a qudit shared between pairs, the act of distilling, a read passed along a pair to its far end and a record fixed by all the reads taken jointly, are the objects every chapter builds with, and the figure beneath draws them once in the conventions the compendium uses throughout. A qudit is a site, a pair is a link between two sites, a shared qudit is a site at two or more links, and a read is an arrow along a link from the qudit read to its far end.
\begin{figure}[htbp]
\centering
\setlength{\unitlength}{1mm}
\begin{picture}(156,36)(0,0)
\put(8,22){\circle{4}}\put(40,22){\circle{4}}\put(10,22){\line(1,0){28}}
\put(8,29){\makebox(0,0){\footnotesize a qudit}}\put(40,29){\makebox(0,0){\footnotesize its far end}}\put(24,18){\makebox(0,0){\footnotesize a pair, one link}}
\put(24,9){\makebox(0,0){\footnotesize the genre of the link, its two end dimensions}}
\put(60,22){\circle{4}}\put(80,22){\circle*{4}}\put(100,22){\circle{4}}\put(62,22){\line(1,0){16}}\put(82,22){\line(1,0){16}}
\put(80,29){\makebox(0,0){\footnotesize a shared qudit, one end of two pairs}}
\put(80,9){\makebox(0,0){\footnotesize a web, pairs joined at shared qudits}}
\put(120,22){\circle{4}}\put(150,22){\circle{4}}\put(122,22){\vector(1,0){26}}
\put(120,29){\makebox(0,0){\footnotesize the qudit read}}\put(150,29){\makebox(0,0){\footnotesize its far end}}
\put(135,18){\makebox(0,0){\footnotesize a read, the basis passed along}}
\put(135,9){\makebox(0,0){\footnotesize distilling, the projection onto a pair}}
\end{picture}
\caption{The objects of the axiom in the conventions of this compendium. A qudit is a site. A maximally entangled pair of qudits is a link between two sites, its genre the pair of its end dimensions. A shared qudit is a site at two or more links, and a web is a set of pairs joined at shared qudits. Distilling two qudits projects them onto a pair on their common support. A read gives a qudit one basis and passes it along each of its pairs to the far end, drawn as an arrow along the link.}
\label{fig:axiom-objects}
\end{figure}
Each chapter of this compendium has one shape. Every result it establishes stands in a numbered box, a theorem, a definition, an open result or a table, in the order of its derivation. Beneath the box one line gives its standing, a theorem, a theorem under stated conditions, an open result or an instrument, and one line names the chapters it rests on and the chapters that rest on it, generated mechanically from the ancestry of the chapters. A short paragraph reads the box in plain words. A paragraph headed in the language of the literature then names what the literature calls each object of the chapter and how far the name reaches, the same object, the nearest one, or a correspondence by name. The derivation follows, its kind named at its head. An exact computation is printed step by step, every intermediate quantity an exact element of the ring in the basis the derivations use, the numbers one, root two, root three and root six, with its smallest instance on the page and its displays numbered. A finite computation is stated as a procedure a reader implements from the text, its smallest instance worked and its result certified at the repository. An exhaustive census of a set gives its counts by kind and names its first parting element, certified likewise. A table is printed in full with the rule that produced it. The derivation ends with a check paragraph that says what was altered and what moved, and, where the result is a finite computation or a census, with the one certificate sentence, the same sentence everywhere, that points at the entry of this chapter at the reproducibility appendix. A chapter whose object is small draws it once, in the conventions of the figure above. Where the reads of the theory admit two readings of a quantity, both are printed side by side and neither is chosen. The reproducibility appendix at the back of the compendium carries one entry per chapter with the names of its checked declarations and the tables its derivation reads, and the Lean package of the compendium is named there as the reproducibility bundle, a package that certifies the statements of this compendium under the definitions of this compendium. An observed value enters the compendium only as a display beside a derived figure, with its provenance. The compendium consumes no measurement. The chapters are gathered in fourteen Parts, whose map stands in the front matter, and a chapter whose result has not yet been derived carries a short note in place of a text.
\begin{chbridge}
The nearest standard term of every word of this compendium stands at the glossary at the front of the compendium, each word tagged as the same object (exact), as a neighbor (nearest), or as a correspondence by name.
\end{chbridge}
Nothing stands open at the framework itself; the notation table of the derivations stands at the repository, and the glossary at the front of this compendium.
\section{The axiom and the six definitions}\label{mod:I-1-axiom-and-the-six-definitions}
\chdesig{1-1}{}{I-1-axiom-and-the-six-definitions}{I-1-D541156-wslv-t1}
\label{narr:open:1-1}This chapter begins the thread at the sentence itself. Its results test how much of the axiom and the five definitions is forced by the rest: Theorem~2.1 obtains the fourth definition from the axiom, Theorem~2.3 obtains one act from the derived act, and Theorem~2.9 excludes the one-dimensional end. The open results Open result~2.4, Open result~2.5, Open result~2.7 and Open result~2.8 name the clauses whose discharge still parts, Theorem~3.1 fixes the far end direction by direction as the read, Open result~4.1 locates the first atom at which the minimality of the text parts, and the two census tables close the chapter.
\begin{chtheorem}[The fourth definition from the axiom]\label{box:I-1-axiom-and-the-six-definitions:1}
This chapter says which of the six definitions are corollaries of the Axiom, which clauses remain premises, and which words are presentations. The definition \emph{known} is a corollary of the Axiom at both clauses: what is known of a qudit is what all its reads (the read: the basis passed along each pair to its far end) fix jointly with certainty whichever way the basis directions are matched across its pairs, and two reads commuting only up to a sign are a further record with their order as the sign, by the holonomy sign at dimension two.
\end{chtheorem}
\chstanding{theorem}
\chdeps{Definitions 1.1 to 1.6}{the theorems of Chapters 1-1 to 1-6, 3-1 to 3-3, 5-1, 5-7, 7-1, 15-1, 15-4, 15-6}
\begin{chconditional}[Orbit reading from the knowledge clause]\label{box:I-1-axiom-and-the-six-definitions:2}
Conditional on the results this theorem rests on at their own standing, as its derivation names them.
The quantifier of \emph{known} binds what is fixed, not the reads or the records: the one attachment leaving every clause of the text non-vacuous fixes the state of the two copies of the pair, invariant under every matching.
\end{chconditional}
\chstanding{conditional theorem}
\chdeps{Definitions 1.1 to 1.6}{the theorems of Chapters 1-1 to 1-6, 3-1 to 3-3, 5-1, 5-7, 7-1, 15-1, 15-4, 15-6}
\begin{chtheorem}[One act from the derived act]\label{box:I-1-axiom-and-the-six-definitions:3}
The one-act clause is a theorem of the derived act, the rank-one operator between pair states (the pair ket) on the common support at every link type (the genre): every entry of the census of moves then carries at most one act, a subdivision being a path of four one-pair acts.
\end{chtheorem}
\chstanding{theorem}
\chdeps{Definitions 1.1 to 1.6}{the theorems of Chapters 1-1 to 1-6, 3-1 to 3-3, 5-1, 5-7, 7-1, 15-1, 15-4, 15-6}
\begin{chwall}[One act clause discharge]\label{box:I-1-axiom-and-the-six-definitions:4}
Read from the text alone the clause is a premise, the text naming no move; under the grammar of one read or one join per link it holds at the single-move kinds and fails first at the subdivision.
\end{chwall}
\chstanding{open result, a derivation carried to a named step at which it parts and printed with its figures on both sides}
\chdeps{Definitions 1.1 to 1.6}{the theorems of Chapters 1-1 to 1-6, 3-1 to 3-3, 5-1, 5-7, 7-1, 15-1, 15-4, 15-6}
\begin{chwall}[Record clause from the one act read]\label{box:I-1-axiom-and-the-six-definitions:5}
The structural half of the clause on records is not forced: the two slots of a record of the one link lie at different ends at ninety-six records and at one end at seventy-six, failing first at the record of two reads at the centre.
\end{chwall}
\chstanding{open result}
\chdeps{Definitions 1.1 to 1.6}{the theorems of Chapters 1-1 to 1-6, 3-1 to 3-3, 5-1, 5-7, 7-1, 15-1, 15-4, 15-6}
\begin{chtheorem}[The restated text at both faces at every cell]\label{box:I-1-axiom-and-the-six-definitions:6}
The restated text is a presentation of the axiom text: all thirty-two derived cells with a phase input read the same at both of its faces, and the re-spelling of the tenth from the ninth is an instrument, the ninth remaining the text in force.
\end{chtheorem}
\chstanding{theorem}
\chdeps{Definitions 1.1 to 1.6}{the theorems of Chapters 1-1 to 1-6, 3-1 to 3-3, 5-1, 5-7, 7-1, 15-1, 15-4, 15-6}
\begin{chwall}[The act alphabet]\label{box:I-1-axiom-and-the-six-definitions:7}
The two letters of the act alphabet, attach at an open far end and join at a closed end, spell the pairing move and the body cut but not the hub step, parting first at the kept leaf of the electron; the carrier records attach at one site each at all but eight of the photon-web records, the exception the second made link of the muon.
\end{chwall}
\chstanding{open result}
\chdeps{Definitions 1.1 to 1.6}{the theorems of Chapters 1-1 to 1-6, 3-1 to 3-3, 5-1, 5-7, 7-1, 15-1, 15-4, 15-6}
\begin{chwall}[The unknowability clause at the axiom]\label{box:I-1-axiom-and-the-six-definitions:8}
The unknowability clause, that the whole is known only by distilling (the distilling act: the projection onto a pair on the common support), is forced at neither reading.
\end{chwall}
\chstanding{open result}
\chdeps{Definitions 1.1 to 1.6}{the theorems of Chapters 1-1 to 1-6, 3-1 to 3-3, 5-1, 5-7, 7-1, 15-1, 15-4, 15-6}
\begin{chtheorem}[The one-dimensional end excluded at the axiom]\label{box:I-1-axiom-and-the-six-definitions:9}
A pair with an end of dimension one is a product state, so the word maximally entangled, read at its predicate, excludes the one-dimensional end; the reading on which the word is a label is printed beside it.
\end{chtheorem}
\chstanding{theorem}
\chdeps{Definitions 1.1 to 1.6}{the theorems of Chapters 1-1 to 1-6, 3-1 to 3-3, 5-1, 5-7, 7-1, 15-1, 15-4, 15-6}
The first clause of \emph{known} follows from the Axiom with the definitions of distilling, read and shared and the standard facts about the pair state, the transport of a basis along a pair and the fixed algebra; the second follows by the Weyl relation, since at a qubit end the commutator phase of two reads is a sign and their product one further traceless operator fixed under every matching, while a qutrit end carries a cube root. The quantifier is located by parsing the clause into its two sentences and testing the three attachments on the record space of the one link. The adopted act is recomputed at every link type and the census of moves re-typed under it. The thirty-two phase cells are read at both faces of the restated text. The Schmidt rank of the flat pair is the smaller dimension at every case tried.

\begin{chbridge}
Two generalized Pauli operators of a $d$-level system obey the Weyl relation $XZ = \omega ZX$ with $\omega$ a primitive $d$-th root of unity: at a qubit they commute up to a sign, at a qutrit up to a cube root \cite{Weyl1927}. The derivation supplies the act, the read and the records of the one link as objects of the pair ket here, and the stabilizer invariant and the swap are named by correspondence, the Born count standing at chapter 1-3 as a dressed measure.
\end{chbridge}
\begin{chproof}[census]
The premises are the axiom and its six definitions (chapter 1-1, the text), the one-act clause as a theorem of the derived act (chapter 1-1, a theorem) and the listed webs with their end lines (chapter 1-5, a theorem). The chain back to the axiom runs through chapter 1-1 and the framework chapter. By exhaustive enumeration of the three thousand three hundred eighty-two link copies of the listed webs (the set generated by every link of every one of the sixty-four listed webs at its two orientations, each end typed open, closed and equal or closed and washed by the definition of open, the far end open when the near dimension exceeds the far). The property tested being whether the two letters of the act alphabet, attach at an open far end and join at a closed end of equal dimension, spell each of the fifty-nine derived moves, the hub step at three bodies, the pairing at the deuteron and the alpha, the body cut at both and the dark-pair move at the forty-nine body rows. The counts by kind, at the pair orientation the ends open at nine hundred thirty, closed and equal at one thousand five hundred twenty-two and closed and washed at nine hundred thirty, with attach admitted under the valence bound at two hundred twenty-one ends and join at thirty-six, neither at nine hundred thirty and an attach or join refused at a saturated site at two thousand one hundred ninety-five. The moves spelled with the product equal at four of fifty-nine at the reading taken, the two pairings and their two cuts,
\begin{equation}
\text{spelled by the two letters: } 4 \text{ of } 59,\qquad \text{at exactly two of the sixteen readings: } 59 \text{ of } 59 ,
\end{equation}
the first parting element the kept leaf of the electron, a (2,4) link at the hub whose ends are closed and washed, the missing letter an attach at a washed end, so the hub step is a third letter, which is the open result. The list end lines matched one to one and the pairing rule re-met at every copy. The smallest case on the page is the deuteron, its one join at the two qutrit ends of its first sites is the pairing move, equal to its listed block, and the reverse of that join is its body cut.

The counter-tests swapped the letters at every end, which moved two hundred thirty-eight cells beginning at the first parting. typed the open ends by the list end lines, which moved one hundred twenty-six. spelled the hub step with the join alone, which moved one hundred eighty-six. paired at free valence zero, which moved one hundred sixty, and struck the text at its gate. Every control read the same at both engines, and the letters at every copy were reproduced at a second engine, equal at every cell.
\chcert{3382}
\end{chproof}
\begin{chtheorem}[The single read, a corollary of the definitions]\label{box:I-1-axiom-and-the-six-definitions:10}
No second copy of a pair exists, so a read acts once on the one copy of each pair the qudit meets, retains the outcome and the basis passed to the far end, and has no statistics over fresh copies of the pair to draw on.
\end{chtheorem}
\chstanding{corollary of the axiom and the definition of a read, derived on this page and checked by no declaration of the Lean package at this release}
\chdeps{Definitions 1.1 to 1.6}{the theorems of Chapters 1-1 to 1-6, 3-1 to 3-3, 5-1, 5-7, 7-1, 15-1, 15-4, 15-6}
The corollary reads the clause \emph{finitely many} of the axiom together with the definition of a read, and nothing else. It names the single-copy task the theory sets every read: the pairs that exist are a finite set and nothing exists beside them, so a pair is never prepared again for a second read, and what a read fixes it fixes at the one copy. This is a use of the clause \emph{finitely many}, which the minimality question of this chapter found consumed in substance by no theorem of the package and by no item of the compendium at both of its faces.
\begin{chbridge}
The literature calls the setting the single-copy regime of quantum state identification (nearest): one copy of a state admits one measurement, no cloning supplies a second copy \cite{WoottersZurek1982}, and state tomography, which estimates a state from the statistics of many fresh copies, is not available. The correspondence is by name, the one copy of a pair following from the finite count of the axiom and not from any dynamics.
\end{chbridge}
\begin{chproof}[exact]
\textit{The premises.} The axiom of Section~\ref{mod:FRAMEWORK}, a definition, at its clauses \emph{nothing exists but} and \emph{finitely many}; the definition of a read of Section~\ref{mod:FRAMEWORK}, a definition, at its clauses \emph{one read for each basis}, \emph{its only basis} and \emph{passes that basis along each pair to its far end}. Nothing outside that list is used, and the chain from the axiom to this result is one chapter deep.

\textit{The steps.} From the clause \emph{finitely many}, by reading it as a count, we obtain that the pairs are a finite set, each pair a member of it once. From the clause \emph{nothing exists but}, by reading it as a closure, we obtain that this set is everything there is, so a second copy of a pair, which would be a further object, does not exist, and nothing exists that could prepare one. From the definition of a read and the first two steps, by reading a read as a distillation of a qudit with another in one of its bases, we obtain that a read acts on the one copy of each pair the qudit meets and, by the clause \emph{one read for each basis}, acts on it once at that basis. From the clause \emph{its only basis} and the clause \emph{passes that basis along each pair to its far end}, by the second step, we obtain that the basis the read gives the qudit, its outcome, and the basis it passes to the far end are retained at the one copy of the pair and at no other object, since no other object carries that pair. From the four steps, by reading a statistic over fresh copies as a count over second copies of the same pair, we obtain that no read has such a statistic to draw on, which is the result.

\textit{The smallest instance.} One pair of two qubits, the one copy of the (2,2) link; a read of one end in one basis passes that basis along the link to the other end; a second read of the same end in the same basis is the same read, one read for each basis; a further qubit pair of the set, where the axiom holds one, is another pair and not a copy of this one.

\textit{The check.} The clause \emph{finitely many} struck, so that pairs are supplied without bound, lets a fresh copy of a pair be read anew and the statistics over copies return, which breaks the fifth step. The clause \emph{one read for each basis} struck lets a qudit be read twice in one basis, which breaks the third step. The result was re-derived by a second route, the two definitions parsed at their clauses by hand with no step beyond them.
\end{chproof}
\section{The far end fixed direction by direction as the read}\label{mod:I-1-D535360-ndfe-t1}
\begin{chtheorem}[The far end fixed direction by direction as the read]\label{box:I-1-D535360-ndfe-t1:1}
At the adopted read on a qubit-qutrit link, where the near qutrit gives its only basis at a cap read and the basis is passed along the pair to the far qubit, the far end is fixed and freed direction by direction by the basis the read passes: every near basis direction whose transport is nonnull fixes exactly one far direction, the traceless part of its projector, and frees the two transverse to it, at all eleven nonnull directions, while the one null direction frees the far end. The washed read along $Z$ fixes the one far letter $Z$ and frees $X$ and $Y$, and the three open reads $X$, $XZ$ and $XZ^{2}$ fix none and free three, their far directions a trine; the fixing is the same at every move of the identification group on the support, and the far marginal is invariant on the state at every read. On this reading the premise that the far end is fixed and freed direction by direction is a tautology of the definition of transport to the far end, drafted as Lean text with its Lean proof still to be written.
\end{chtheorem}
\chstanding{theorem}
\chdeps{Definitions 1.1 to 1.6 and Theorems 2.1 to 2.3, Open results 2.4 and 2.5, Theorem 2.6, Open results 2.7 and 2.8 and Theorems 2.9 and 2.10}{no theorem beyond it at this release}
The adopted read at the one qubit-qutrit link is taken direction by direction at the near qutrit; for every near basis direction the transport along the pair is computed and its far image read as the traceless part of a projector, a unit element of the derived direction triple; the fixed and freed far directions are counted at every nonnull direction and at the one null direction, the invariance under the identification group is checked over its monomials, and the far marginal is compared across the reads. The result is a theorem by construction and names no condition beyond the results it rests on.

\begin{chbridge}
A measurement on one member of an entangled pair fixes the conditional state of the other member in the measured basis and leaves the complementary directions undetermined, the steering of Schrodinger and the correlations of Einstein, Podolsky and Rosen \cite{Schrodinger1935,EPR1935}; the far marginal is unchanged by the choice of basis, which is the no-signaling condition. The derivation supplies the transport of a basis to the far end and the far directions it fixes and frees here, and the helicity and the handedness it would label are named by correspondence.
\end{chbridge}
\begin{chproof}[finite and census]
The premises are the pair ket and the transport to the far end (chapter 1-1, a theorem), the link with its six types (chapter 1-2, a theorem) and the record phase group at every type (chapter 1-2, a theorem). The chain back to the axiom runs through chapter 1-1 and the framework chapter. The procedure. The input is the one (2,3) link with the near qutrit, the far qubit and their common support of dimension two, and the four cap bases of the near end, the washed read $Z$ and the three open reads $X$, $XZ$ and $XZ^{2}$, twelve near basis directions in all. For every near direction the transport along the pair is computed as the support projection of the conjugate, giving a far vector as an exact element of the ring in the basis one, root two, root three, root six. The test applied to each is whether the far vector is nonnull, in which case the far direction it fixes is read as the traceless part of its projector, a unit element of the derived direction triple, with the two directions transverse to it free, or null, in which case the far end is freed. The aggregation counts the nonnull directions fixing exactly one far direction, the null directions, the fixed and freed far letters per read, the covariance of the fixing at every monomial of the identification group and the invariance of the far marginal. The smallest instance on the page is the washed read $Z$ at the near qutrit, its three directions transport to the far vectors
\begin{equation}
\lvert 0\rangle \mapsto (1,\,0),\qquad \lvert 1\rangle \mapsto (0,\,1),\qquad \lvert 2\rangle \mapsto 0 ,
\end{equation}
the first fixing the far letter $Z$ at the unit $+1$, the second fixing the same letter at $-1$, the third null and freeing the far end, so the washed read fixes one far letter and frees the two transverse, $X$ and $Y$. The open read $X$ transports its three directions to the far vectors $(1,1)$, $(1,-1+\sqrt{3})$ and $(1,-\sqrt{3})$, three nonnull vectors fixing no common letter and freeing three, a trine. The totals at the checkpoints,
\begin{equation} \adjustbox{max width=\dimexpr\linewidth-3em\relax}{$\displaystyle \begin{gathered} \text{nonnull directions fixing exactly one far direction: } 11 \text{ of } 11, \qquad \text{null directions: } 1, \\ \text{covariance: } 384 \text{ of } 384, \qquad \text{far marginal invariant: } 48 \text{ of } 48 , \end{gathered}$} \end{equation}
the covariance taken over the thirty-two monomials of the identification group with fourth-root phases at the twelve directions, and the result, the far end is fixed and freed direction by direction by the basis the read passes along the pair, so the premise that it is so is a tautology of the definition of transport to the far end, which is the theorem. The same count at the dock face of the (2,3) type, six of six nonnull directions and the washed reads at one fixed and two freed, stands beside.

The counter-tests struck the pair ket, which moved twelve of thirteen cells beginning at the transport. struck the conjugate, which moved seven. struck the near basis, which moved nine. struck the fixing rule, which moved the count cells, and struck the text at its gate, which moved every cell. Every control read the same at both engines, and the fixed and freed directions were reproduced at a second engine, equal at every cell.
\chcert{384}
\end{chproof}
\section{The minimality question for the six definitions}\label{mod:I-1-D535761-w-mnth}
\begin{chwall}[The minimality of the text at its first atom]\label{box:I-1-D535761-w-mnth:1}
The minimality of the axiom text parts at its first atom. The words finitely many of the axiom are consumed in substance by no theorem of the Lean package and by no item of the compendium, every finiteness token of the items standing at the matching group, so the finiteness is the first word of the axiom the derivation does not use; beside it, the derivation graph of the definitions is cyclic at the adopted reading, the second clause of \emph{read} citing what is known and the first clause of \emph{known} citing the definition of a read. The opposite face, that every atom of the text is needed with its consumer named and derived by no other, holds at the six axiom atoms and at the five definition atoms whose removal breaks the build, with the gauge atom consumed by no readout but moving the phase residue; the one atom typed unused at both faces is the finiteness atom.
\end{chwall}
\chstanding{open result}
\chdeps{Definitions 1.1 to 1.6 and Theorems 2.1 to 2.3, Open results 2.4 and 2.5, Theorem 2.6, Open results 2.7 and 2.8 and Theorems 2.9 and 2.10}{no theorem beyond it at this release}
The census of the atoms of the axiom text is re-met, the consumer of every atom is named and quoted at the Lean package and at the compendium items, the derivation graph is built at two readings with its faces and edges, the finiteness tokens are counted at every item, and the gauge atom is tested by rephasing the Gram elements; the result stands open because the finiteness words are consumed by no declaration and by no item in substance.

\begin{chbridge}
An axiom system is minimal when every axiom is independent of the others, each one needed for some theorem and derivable from none of the rest; the independence proofs of Hilbert for the axioms of geometry are the model, and the finiteness assumptions of a quantum axiom system, a finite-dimensional state space, are among the assumptions whose use is traced theorem by theorem \cite{Hilbert1899,Mackey1963}. The correspondence is by name.
\end{chbridge}
\begin{chproof}[exact]
The premises are the axiom and its six definitions (chapter 1-1, the text), the pair ket and the transport to the far end (chapter 1-1, a theorem) and the readability order at the six types (chapter 1-2, a theorem). The chain back to the axiom runs through chapter 1-1 and the framework chapter. From the axiom text, by the census of its atoms (every word or clause an item of the compendium or a declaration of the package consumes), we obtain nineteen atoms typed at the face that holds, six atoms of the axiom needed with their consumer named and derived by no other, twelve atoms of the definitions that are corollaries with their derivation, and one gauge atom, the phase residue consumed by no readout,
\begin{equation}
6 + 12 + 1 = 19,\qquad \text{atoms unused at the census} = 0 .
\end{equation}
From the items of the compendium and the declarations of the package, by the count of the tokens that consume the words finitely many of the axiom, we obtain the consumption cells of the finiteness atom, zero pair-set finiteness tokens at the distilling corollary, every finiteness token of the items standing at the matching group, the exhaustion words nothing exists but consumed at one item, and the finiteness consumed at the word by zero items,
\begin{equation}
n_{\text{finiteness, in substance}} = 0,\qquad n_{\text{finiteness, at the word}} = 0,\qquad n_{\text{exhaustion}} = 1 ,
\end{equation}
so the finiteness atom is unused at both faces, the first atom of the axiom the derivation does not consume, and the open result stands there. From the definitions with their clauses, by the derivation graph at the proof-rooted reading (an edge from a clause to the definition its proof cites), we obtain a graph that is cyclic at one component, the second clause of \emph{read} cites what is known and the first clause of \emph{known} cites the definition of a read. At the literal reading the cycle widens to five clauses of the two definitions. From the gauge atom, by the rephasing of the letters at dimension two, we obtain two hundred eighty-eight moved elements of the Gram matrix and no moved modulus, so the gauge atom is consumed by no readout and moves the phase residue alone. The opposite face, that every atom is needed with its consumer named and derived by no other, holds at the six atoms of the axiom and at the five atoms of the definitions whose removal breaks the build of the package, re-met equal to the census at the strike of each declaration. The smallest instance on the page is the cycle of two clauses, the clause of \emph{read} that cites what is known and the clause of \emph{known} that cites a read, each the consumer of the other.

The counter-tests struck each of the five definition atoms that carry a declaration, which broke the build of the package at every one, altered the axiom to an infinite set, which moved no consumption cell, and rephased the gauge atom, which moved the Gram elements and no modulus. The controls stood at their pinned values but one, which the derivation disclosed with the result. The census of the atoms was re-met at a second route, equal at every atom.
\end{chproof}
\begin{figure}[h]\centering
\begin{adjustbox}{max width=\linewidth}\begin{tabular}{lrl}\toprule
atoms of the axiom text & count & the face that holds \\ \midrule
atoms of the axiom needed and derived by no other & 6 & each consumed by a named item \\
atoms of the definitions that are corollaries & 12 & each with its derivation \\
the gauge atom & 1 & the phase residue, consumed by no readout \\
the finiteness atom among the six & 1 & consumed in substance by nothing \\ \bottomrule
\end{tabular}\end{adjustbox}
\caption{The atoms of the axiom text as the census types them, nineteen in all, with the finiteness atom the one unused at both faces.}
\end{figure}
\section{The definition-choice invariance census, first count}\label{mod:I-1-D538240-the-definition-choice-invariance-census}
\begin{chtable}[The definition-choice invariance census, first pass]\label{box:I-1-D538240-the-definition-choice-invariance-census:1}
This instrument computes the pinned cells at both faces of each of the three definition choices of the derivation, the finiteness clause of the axiom, the stop condition of the cap census and the adopted reading of the definition of open, at eight faces over eighty-five cells pinned ahead and recomputed at every face by the derived engines, and types each choice by a pinned rule. The definition of open is typed a premise at its adopted reading, its first parting cell the capped seed; the stop condition is left incomplete at a cell by a reader fault of the census itself, cured at its second pass; and the join reading is typed a premise at its adopted reading, its first parting cell the pairable star.
\end{chtable}
\chstanding{instrument}
\chdeps{Definitions 1.1 to 1.6 and Theorems 2.1 to 2.3, Open results 2.4 and 2.5, Theorem 2.6, Open results 2.7 and 2.8 and Theorems 2.9 and 2.10}{Table 6.1}
Every cell of the three choices is recomputed by the derived engines at every face (the depth result at its four readings of open, the bound-respecting result at its letters, the body-alone and count results at the five bodies, the lattice, many-body and nuclear results at the join); the typing rule is applied to the cells and types each choice from its cells; the first parting cell of every premise is named.

\begin{chbridge}
A physical theory stated from an axiom and a set of definitions must be checked for its dependence on the choices made at the definitions: a prediction that moves when an equivalent definition is substituted is a convention and not a result, and the invariance of the predictions under the admissible choices is the test that the definitions are well posed \cite{Einstein1905,Wigner1960}. The derivation supplies the cells of the three choices at their faces here, and the moduli are named by correspondence.
\end{chbridge}
\begin{chproof}[instrument]
The table of this instrument has one row per definition choice of the derivation, three rows, and a column at each of the eight faces of the choice, the faces being the readings of the definition of open (the key reading, the shared reading, the valence reading and the key-and-shared reading) and the letters of the stop condition and of the join. The cells are the eighty-five pinned cells of the derived results recomputed at every face by the derived engines, the depth result at its four readings of open, the bound-respecting result at its letters, the body-alone and count results at the five bodies, and the lattice, many-body and nuclear results at the join. The rule that produced the typing, a choice is a convention when every compared cell is equal across its faces, a premise when a compared cell parts, with the first parting cell named and the choice owed to a further derivation, and incomplete at a cell when a compared cell is absent at a face. The rows as they fell:
\begin{equation} \adjustbox{max width=\dimexpr\linewidth-3em\relax}{$\displaystyle \begin{array}{llll}
\text{choice} & \text{face taken} & \text{typing} & \text{first parting cell}\\
\text{the definition of open} & \text{the key reading} & \text{a premise} & \text{the capped seed: } 221,\ 497,\ 485,\ 214 \text{ at the four readings}\\
\text{the stop condition} & \text{the second letter} & \text{incomplete at a cell} & \text{the seed letter absent at the reader of this pass}\\
\text{the join} & \text{the dock reading} & \text{a premise} & \text{the pairable star: three sites against one}
\end{array}$} \end{equation}
the typing two of three, the stop condition left incomplete by a reader of this pass that split a capture line at a tab alone where three entries print a space, a fault of the pass itself found beyond its pin and cured at the second pass. The compared census of the stop condition beside at five cells equal and six parting, the first parting the sinks at the first wave, nine against eleven.

The counter-tests compared the gauge moduli, which read a convention as pinned. typed the definition of open at the key faces, a convention as pinned. replaced the letters of the stop condition at a copy, which moved no cell where a move was pinned, the reader fault. struck the text at its gate, which moved every cell, and typed the definition of open incomplete at a copy with the capped seed absent, as pinned. Every control read the same at both engines, eight of eight, and the cells were reproduced at a second engine, equal at every face.
\end{chproof}
\section{The definition-choice invariance census, second count}\label{mod:I-1-D538500-the-definition-choice-invariance-census}
\begin{chtable}[The definition-choice invariance census, second pass]\label{box:I-1-D538500-the-definition-choice-invariance-census:1}
The second pass of the census re-fires the first with the first as its base, the stop condition typed at a reader that reads a capture line at its first tab or space, so that the three definition choices of the derivation stand typed by the pinned rule at every face: the definition of open a premise at its adopted reading with its first parting cell the capped seed, the stop condition of the cap census a premise at its adopted reading with its first parting cell the first-wave sinks, and the join reading a premise at its adopted reading with its first parting cell the pairable star; three choices typed, the census complete.
\end{chtable}
\chstanding{instrument}
\chdeps{Definitions 1.1 to 1.6, Theorems 2.1 to 2.3, Open results 2.4 and 2.5, Theorem 2.6, Open results 2.7 and 2.8 and Theorems 2.9 and 2.10 and Table 5.1}{no theorem beyond it at this release}
The eighty-five pinned cells are recomputed at every one of the eight faces by the derived engines, the reader fault of the first pass cured ahead of the pin, and the typing rule is applied to the cells, typing each choice from its cells alone.

\begin{chbridge}
The invariance of the predictions of a theory under the admissible choices at its definitions is the test that the definitions are well posed: a quantity that moves when an equivalent definition is substituted is a convention and not a result, and a census of the choices at every face is the systematic form of that test, the counterpart of the gauge and coordinate independence demanded of a physical quantity \cite{Einstein1905,Wigner1960}. The derivation supplies the cells of the three choices at their faces here, and the moduli are named by correspondence.
\end{chbridge}
\begin{chproof}[instrument]
The table of this second pass has the same three rows and eight faces as the first, the eighty-five pinned cells recomputed at every face by the derived engines with the first pass as its base, the reader of this pass reading a capture line at its first tab or space so that the cells the first pass left absent are read. The rule that produced the typing is the rule of the first pass, a convention when every compared cell is equal across the faces, a premise when a compared cell parts with the first parting cell named and the choice owed to a further derivation, and incomplete at a cell when a compared cell is absent. The rows as they fell:
\begin{equation} \adjustbox{max width=\dimexpr\linewidth-3em\relax}{$\displaystyle \begin{array}{llll}
\text{choice} & \text{face taken} & \text{typing} & \text{first parting cell}\\
\text{the definition of open} & \text{the key reading} & \text{a premise} & \text{the capped seed: } 221,\ 497,\ 485,\ 214 \text{ at the four readings}\\
\text{the stop condition} & \text{the second letter} & \text{a premise} & \text{the sinks at the first wave: } 9 \text{ against } 11\\
\text{the join} & \text{the dock reading} & \text{a premise} & \text{the pairable star: three sites against one}
\end{array}$} \end{equation}
the typing three of three with no choice untyped, no convention and no cell incomplete.

The counter-tests compared the gauge moduli, which read a convention as pinned. typed the definition of open at the key faces, a convention as pinned. replaced the stop-condition cells at a copy by one value at both faces, which moved the typing to a premise at other cells as pinned. struck the text at its gate, which moved every cell, and typed the definition of open incomplete at a copy with the capped seed absent, as pinned. Every control read the same at both engines, and the cells were reproduced at a second engine, equal at every face.
\end{chproof}
\section{The definition of open at the larger end}\label{mod:I-1-D539391-w-dole}
A pure state of a bipartite system has a Schmidt decomposition, a sum of product terms over orthonormal bases of the two factors with non-negative weights, and the number of nonzero weights (the Schmidt rank) is bounded by the smaller of the two local dimensions; for a maximally entangled pair of unequal local dimensions the larger factor carries a subspace that no term of the decomposition reaches \cite{Schmidt1907,Nielsen2000}. The terms used below are introduced in the framework chapter (section~\ref{mod:FRAMEWORK}), the axiom and the six definitions at section~\ref{mod:I-1-axiom-and-the-six-definitions} and the link with its six genres at section~\ref{mod:I-2-link-and-its-six-genres}.
Take as the definition of open the identity that the three-premises result (section~\ref{mod:II-9-D539176-w-tprd}) found at every end: an end of a link is open when its dimension exceeds the dimension of the far end, so that the larger end of every unequal link is open. Under this derived definition the open set coincides, at every one of the 1262 ends of the eight webs of the cap census, with the set of ends the read leaves open at the far end and with the set of ends the wash still moves on the state, at 1262 of 1262 ends in each case.
The derivation re-reads every end of the 631 links of the eight webs, the universe web and the seven roster webs, under the derived definition and compares three sets end by end: the derived open set, the read-open set (the ends whose passed directions fall short of their own dimension) and the wash-moved set; the three coincide at every end, so the identity is reproduced. The derived definition is the mirror image of the standing one, which marks the smaller end of an unequal link open: at every end the standing letter equals the derived letter at the far end, the mirror holding at 1262 of 1262 ends, and because neither definition marks an end of an equal link open, the two open sets agree end by end at 500 of the 1262 ends and part, by construction, at the two ends of every unequal link. The results that consume the standing definition were then recomputed under the derived one: the cap census reproduces 6 of its 11 values, the two port results 14 of 22 and 49 of 85, the horizon metric 6 of 10 and the stop condition 25 of 25, so that one of the five is reproduced whole; the first value to move is the census of the uncapped end sets of the cap census, whose 330 uncapped links are read at the centre end under one definition and at the far end under the other.
The result is an open one, an open result in the terms of this compendium: the derived definition reproduces the identity at every end, but the results that consume the definition are not all reproduced under it, and the first to part is named above.
What stays open is whether a reading of open exists under which every consuming result is reproduced; the link typing at the cap (section~\ref{mod:I-2-D539697-w-aptc}) and the dressing as the Higgs-to-matter leg (section~\ref{mod:VI-10-D539753-w-yklg}) read the derived definition at their own questions.
\section{The washed level as an output of the definition of open}\label{mod:I-1-D541156-wslv-t1}
\begin{chtheorem}[The washed level as an output of the definition of open]\label{box:I-1-D541156-wslv-t1:1}
Of the three wash levels, the root, the washed and the flat, as the three readout rules carry them to a state of the records, exactly one is an output of the definition of open at the second-order office of the four words, the level washed: at the key reading applied at the one-act read under the slot rule, the far end of every dark link of the W word, the Z word, the Higgs word and the gluon pair ket is washed, the ququart far ends of the (2,4) legs and the equal ends of the (4,4) and (3,3) links, the level set at the dark links of the four words is the singleton washed at four of four words, and the second-order offices of the Higgs and carrier words are re-met at that level at four of four words, the rate law reading the closure of the body-alone office as the sum of the squared readouts at the washed level at every word, while the gain lemma at the root level and the readout unit at the flat level read it at none. The second face is unmet and the first parting is none.
\end{chtheorem}
\chstanding{theorem}
\chdeps{Definitions 1.1 to 1.6, Theorems 2.1 to 2.3, Open results 2.4 and 2.5, Theorem 2.6, Open results 2.7 and 2.8 and Theorems 2.9 and 2.10, Theorems 18.1 to 18.4 and Open results 18.5 and 18.6, Theorem 195.1 and Theorems 200.1 to 200.3 and Open result 200.4}{no theorem beyond it at this release}
The three levels are read at every word as three states of the records over its dark links, the pure state of the readout vector at the root, the state beyond one act of the wash at the washed level and the fixed point of the wash at the flat level, and at each level the three readout rules are compared with the closure of the body-alone office; the level set at the dark links of each word is typed under the six readings of the definition of open and under the slot, act and tier readings beside, and the face is read at the reading named.

\begin{chbridge}
In the standard formalism a measurement is a projective read followed by a partial trace over the registers that are not observed, and the rate of a process is the squared matrix element of the transition summed over the unobserved final states, the golden rule \cite{Dirac1927,vonNeumann1932}; the coherent sum over amplitudes stands at the root of that hierarchy and the fully mixed state at its end. The derivation supplies the three states and the one level at which the closure is re-met here, and the decoherence and the rate it would fix are named by correspondence.
\end{chbridge}
\begin{chproof}[census and exact]
The premises are the three readout rules at the three levels of the wash (chapter 1-3, a theorem), the wash move at the cap (chapter 11-6, a theorem), the body-alone office of every word by the one rule (chapter 6-10, a theorem), the body-alone square at the five bodies (chapter 8-2, a theorem) and the two definitions of open at every end (chapter 1-1, an open result). The chain back to the axiom runs through chapter 1-3 and the framework chapter. By exhaustive enumeration of the one hundred forty-four readings of the typing, the six readings of the definition of open, the three slot readings, the two act counts, the two tiers and the two gluon readings, at the four words, the property tested at each being whether the level set at the dark links of the four words is a singleton and whether the second-order offices are re-met at that level at four of four words. The counts by kind are the level set a singleton at fifty-six of the one hundred forty-four readings and a singleton with the offices re-met at it at twenty, among them the reading named, the key reading at the slot rule, the one-act read, the dark-links tier and the pair ket, where the level set is the singleton washed and the offices are re-met at four of four. At the shared, the valence and the key-and-shared readings the level set is the singleton root with the offices re-met at none of four, the gain lemma reading the coherent sum and never the closure; at the read-open reading at the far end the level set holds both the root and the washed level; at the cap reading the level set is the singleton flat with the offices re-met at none of four.
The smallest instance on the page is the W word at its two (2,4) legs, the first power twenty-eight. From the act read of each leg and its second-order readout, by the closure rule, we obtain the closure over the two dark links,
\begin{equation}
c_{W} = \tfrac{10241}{3645000} + \tfrac{98}{50625}\sqrt{2} ,
\end{equation}
and from the first power and the closure, by the one rule, the office of the word,
\begin{equation}
\omega_{W} = 28^{2} + c_{W} = \tfrac{2857690241}{3645000} + \tfrac{98}{50625}\sqrt{2} .
\end{equation}
From the state of the records of the two legs, by the three readout rules at their levels, we obtain the three faces at the W word, at the root the coherent sum of the two readouts, $\tfrac{7}{150} + \tfrac{28}{675}\sqrt{2}$, at the washed level the sum of the squared readouts, $\tfrac{10241}{3645000} + \tfrac{98}{50625}\sqrt{2}$, equal to the closure, and at the flat level the unit, one. The same comparison at the Higgs word, the first power sixteen with the closure $\tfrac{1681}{182250000}$ at its one (4,4) link and the office $\tfrac{46656001681}{182250000}$, gives the root face $\tfrac{41}{13500}$, the washed face equal to the closure and the unit; and at the gluon pair ket, the first power nine with the closure $\tfrac{11989}{84934656} + \tfrac{1675}{21233664}\sqrt{3}$ at its one (3,3) link, the root face $\tfrac{67}{9216} + \tfrac{25}{4608}\sqrt{3}$, the washed face equal to the closure and the unit. So the rate law re-meets the closure at the washed level at four of four words, the gain lemma at the root level at none and the readout unit at the flat level at none, which is the theorem.

The counter-tests struck the wash maps to the identity, which moved the three states to one at every word and the offices re-met to two of four with the face to the second. A second reversed the slot reading, which moved the far ends of the (2,4) legs to the qubits and the level set to the root and the washed level together. A third widened the dark set to every link, which moved the (2,3) links into the closure and their qubit far ends into the level set. A fourth struck the text at its gate, which left the gate unmet, and a fifth permuted the level map, which moved the level named to the root with the offices re-met at none. Every control read the same at both engines, sixteen of sixteen, and the typing was reproduced at a second engine, equal at every cell.
\chcert{144}
\end{chproof}
\begin{figure}[h]\centering
\begin{adjustbox}{max width=\linewidth}\begin{tabular}{llll}\toprule
level & readout rule & the face at the W word & equal to the closure \\ \midrule
the root & the gain lemma & $\tfrac{7}{150} + \tfrac{28}{675}\sqrt{2}$ & no \\
the washed & the rate law & $\tfrac{10241}{3645000} + \tfrac{98}{50625}\sqrt{2}$ & yes \\
the flat & the readout unit & $1$ & no \\ \bottomrule
\end{tabular}\end{adjustbox}
\caption{The three levels of the wash at the W word: the face each readout rule reads at its level against the closure of the body-alone office over the two (2,4) legs, re-met at the washed level alone.}
\end{figure}
\par\medskip\noindent{\itshape With the text tightened on itself, the next chapter introduces the object that carries the rest of the compendium, the link and its six genres.}\label{narr:close:1-1}\par
\section{The link and its six genres}\label{mod:I-2-link-and-its-six-genres}
\chdesig{1-2}{}{I-2-link-and-its-six-genres}{I-2-D543154-w-tntp}
\label{narr:open:1-2}The link joins two ends, each with a register of dimension two, three or four, and the six unordered pairs of registers are its six genres. Theorem~9.1 orders the genres by the first order at which a read does not vanish, Theorem~9.2 shows the abelian genre unique, Theorem~9.3 names the coupling genre, and Theorem~9.6 prints the one-link tables at every genre. The chapter adds the readout unit at the one link, Theorem~11.1, the soft limit there, Theorem~12.1, and the record phase group at every genre, Theorem~13.1, with the cap reached and left type by type, Open result~10.1, as an open result.
\begin{chconditional}[Readability order at the six genres]\label{box:I-2-link-and-its-six-genres:1}
Conditional on the results this theorem rests on at their own standing, as its derivation names them.
This chapter types the links. A link is a pair of the Axiom and its type (the genre) is the pair of dimensions of its two ends; at the admitted alphabet, the dimensions two, three and four, the types are six, written (2,2), (2,3), (2,4), (3,3), (3,4) and (4,4), the qutrit-qubit link as (2,3). The first non-vanishing order of the read is defined at every one of the six: (2,2) and (2,3) are readable at order one under the complete-read bound, (3,3) is decided at order two, and so are (2,4), (3,4) and (4,4), a conditional theorem.
\end{chconditional}
\chstanding{conditional theorem}
\chdeps{Definitions 1.1 to 1.6 and Theorems 2.1 to 2.3, Open results 2.4 and 2.5, Theorem 2.6, Open results 2.7 and 2.8 and Theorems 2.9 and 2.10}{the theorems of Chapters 1-2 to 1-5, 3-1, 3-5, 4-6, 5-1 to 5-9, 5-16, 6-1 to 6-7, 6-7a, 7-1, 7-7 to 7-8, 8-2 to 8-3, 9-1 to 9-2, 9-5, 11-1, 12-4, 14-1, 14-3, 15-4, 15-6}
\begin{chtheorem}[Abelian genre uniqueness]\label{box:I-2-link-and-its-six-genres:2}
Exactly one readable type carries an acting centre: the centre of the identification algebra has dimension one at (2,2) and at (2,3), but it acts on the pair space as a scalar only at (2,3), whose open read alone is abelian, conserved, massless and of spin one.
\end{chtheorem}
\chstanding{theorem}
\chdeps{Definitions 1.1 to 1.6 and Theorems 2.1 to 2.3, Open results 2.4 and 2.5, Theorem 2.6, Open results 2.7 and 2.8 and Theorems 2.9 and 2.10}{the theorems of Chapters 1-2 to 1-5, 3-1, 3-5, 4-6, 5-1 to 5-9, 5-16, 6-1 to 6-7, 6-7a, 7-1, 7-7 to 7-8, 8-2 to 8-3, 9-1 to 9-2, 9-5, 11-1, 12-4, 14-1, 14-3, 15-4, 15-6}
\begin{chtheorem}[Coupling genre]\label{box:I-2-link-and-its-six-genres:3}
The coupling office of the chain sum at one link read twice is a function of the link type alone, re-meeting the derived fine-structure cell at (2,3) and printing the other types as exact faces.
\end{chtheorem}
\chstanding{theorem}
\chdeps{Definitions 1.1 to 1.6 and Theorems 2.1 to 2.3, Open results 2.4 and 2.5, Theorem 2.6, Open results 2.7 and 2.8 and Theorems 2.9 and 2.10}{the theorems of Chapters 1-2 to 1-5, 3-1, 3-5, 4-6, 5-1 to 5-9, 5-16, 6-1 to 6-7, 6-7a, 7-1, 7-7 to 7-8, 8-2 to 8-3, 9-1 to 9-2, 9-5, 11-1, 12-4, 14-1, 14-3, 15-4, 15-6}
\begin{chtheorem}[Link lemmas]\label{box:I-2-link-and-its-six-genres:4}
The open (2,3) link carries no mass face, its two-read record at the qubit end carries spin one, and zero register displacement equals zero momentum transfer at the one-link star.
\end{chtheorem}
\chstanding{theorem}
\chdeps{Definitions 1.1 to 1.6 and Theorems 2.1 to 2.3, Open results 2.4 and 2.5, Theorem 2.6, Open results 2.7 and 2.8 and Theorems 2.9 and 2.10}{the theorems of Chapters 1-2 to 1-5, 3-1, 3-5, 4-6, 5-1 to 5-9, 5-16, 6-1 to 6-7, 6-7a, 7-1, 7-7 to 7-8, 8-2 to 8-3, 9-1 to 9-2, 9-5, 11-1, 12-4, 14-1, 14-3, 15-4, 15-6}
\begin{chtheorem}[Ordering theorem]\label{box:I-2-link-and-its-six-genres:5}
The passed-hub count is monotone against the hub dressing (the hub: the dimension-four site of a body; the dressing: the hub dressing of a body) at every pair of bodies and every face.
\end{chtheorem}
\chstanding{theorem}
\chdeps{Definitions 1.1 to 1.6 and Theorems 2.1 to 2.3, Open results 2.4 and 2.5, Theorem 2.6, Open results 2.7 and 2.8 and Theorems 2.9 and 2.10}{the theorems of Chapters 1-2 to 1-5, 3-1, 3-5, 4-6, 5-1 to 5-9, 5-16, 6-1 to 6-7, 6-7a, 7-1, 7-7 to 7-8, 8-2 to 8-3, 9-1 to 9-2, 9-5, 11-1, 12-4, 14-1, 14-3, 15-4, 15-6}
\begin{chtheorem}[The one link tables at every genre]\label{box:I-2-link-and-its-six-genres:6}
The one-link tables are defined at all six types by the machinery of the support-rank result, the (2,3) table agreeing with that result at every cell.
\end{chtheorem}
\chstanding{theorem}
\chdeps{Definitions 1.1 to 1.6 and Theorems 2.1 to 2.3, Open results 2.4 and 2.5, Theorem 2.6, Open results 2.7 and 2.8 and Theorems 2.9 and 2.10}{the theorems of Chapters 1-2 to 1-5, 3-1, 3-5, 4-6, 5-1 to 5-9, 5-16, 6-1 to 6-7, 6-7a, 7-1, 7-7 to 7-8, 8-2 to 8-3, 9-1 to 9-2, 9-5, 11-1, 12-4, 14-1, 14-3, 15-4, 15-6}
\begin{chconditional}[Universe link from parity]\label{box:I-2-link-and-its-six-genres:7}
Conditional on the results this theorem rests on at their own standing, as its derivation names them.
The universe link is forced by parity: the admitted bodies carry an odd number of docks, every docking fills two docks of two bodies, so every maximal docking sequence leaves exactly one dock open, and that dock is the one link of the whole; and the pair state of a universe link of dimension $N$ is, entry by entry, the tensor product of the pair states at the factors of $N$, the readable factors at the valence bound being two, three and four, which recovers the six types.
\end{chconditional}
\chstanding{conditional theorem}
\chdeps{Definitions 1.1 to 1.6 and Theorems 2.1 to 2.3, Open results 2.4 and 2.5, Theorem 2.6, Open results 2.7 and 2.8 and Theorems 2.9 and 2.10}{the theorems of Chapters 1-2 to 1-5, 3-1, 3-5, 4-6, 5-1 to 5-9, 5-16, 6-1 to 6-7, 6-7a, 7-1, 7-7 to 7-8, 8-2 to 8-3, 9-1 to 9-2, 9-5, 11-1, 12-4, 14-1, 14-3, 15-4, 15-6}
\begin{chwall}[Mixed genre functional at the census]\label{box:I-2-link-and-its-six-genres:8}
Three open results stand beside: under the projection matching, the three mixed-type units and the five traceless Cartan weights are conserved at every interaction of the census while the zero set agrees with the bodies named neutral at only nine of thirty deciding pairs; at the fixed list of networks the dimension equals the degree plus the free valence at most sites and parts at the rest, the valence bound broken at one hundred five sites; and with the chain webs set aside the free valence at the act equals the deficit at fewer than half of the retained sites.
\end{chwall}
\chstanding{open result}
\chdeps{Definitions 1.1 to 1.6 and Theorems 2.1 to 2.3, Open results 2.4 and 2.5, Theorem 2.6, Open results 2.7 and 2.8 and Theorems 2.9 and 2.10}{the theorems of Chapters 1-2 to 1-5, 3-1, 3-5, 4-6, 5-1 to 5-9, 5-16, 6-1 to 6-7, 6-7a, 7-1, 7-7 to 7-8, 8-2 to 8-3, 9-1 to 9-2, 9-5, 11-1, 12-4, 14-1, 14-3, 15-4, 15-6}
The readability order is read off the first non-vanishing order of the read at every type, agreeing at both port faces. The center of the identification algebra is found by the adjoint and the Killing routes; at (2,3) its generator survives the twirl at both ends and is central at every matching. The coupling office depends on the two end dimensions alone through the types of records at the one link. The link lemmas follow from the absence of chains at zero displacement, the decomposition of the two-read record under $\mathfrak{su}(2)$ and the vanishing of the signed step term over the star. The parity argument counts docks at every maximal sequence; the tensor-product identity holds at every factorization up to thirty-six, the Schmidt rank being the product of the ranks.

\begin{chbridge}
A pure state of two systems of dimensions $a$ and $b$ is maximally entangled when its reduced operators are flat, its local unitary freedom generated by $\mathfrak{su}(a) \oplus \mathfrak{su}(b) \oplus \mathfrak{u}(1)$ with a one-dimensional centre. The derivation supplies the first non-vanishing order of the read at every type, the center of the identification algebra and the two-point function of the one link here, and the spin, the mass, the momentum transfer and the constant are named by correspondence, the spin as a label of the little group standing at chapters 3-2 to 3-5.
\end{chbridge}
\begin{chproof}[census]
The premises are the axiom with its definition of a read (chapter 1-1, the text), the carrier as the closed read of its type (chapter 4-3, a theorem) and the read engine with its catalog of the one-link reads at order one (an instrument outside the chapters of this compendium). The chain back to the axiom runs through chapter 1-1 and the framework chapter. By exhaustive enumeration of the one thousand six hundred ninety-one links of the listed webs at their six types (the set generated by every link of the sixty-four listed webs, each at the pair of dimensions of its two ends among two, three and four). The property tested at every type being the first non-vanishing order of the read, read at the face taken and at the whole and cut port faces. The counts by kind, the order defined at six of six types and equal at both port faces at six of six,
\begin{equation}
o(2,2)=1,\quad o(2,3)=1,\quad o(2,4)=2,\quad o(3,3)=2,\quad o(3,4)=2,\quad o(4,4)=2 ,
\end{equation}
the readable types (2,2) and (2,3) at order one under the complete-read bound, re-met at the catalog read cell with the reads and the readable set, the (3,3) type at order two re-met at the confinement cells, and the (2,4), (3,4) and (4,4) types decided at order two, the (4,4) type at its closure sector of one hundred five at the gain four with thirty-six sign records at the far ququart end beside. The first undefined type none and the first parting between the faces none, so the readability order is a typing of the type, which is the theorem. At the reading beside, the sign-clause records of the far end as the closure, the orders are defined at four of six with the (3,3) type undefined first, printed beside for comparison. The smallest case on the page is the (2,2) link, its read is readable at order one under the complete-read bound, the first order at which it does not vanish.

The counter-tests of the derivation stood at their pinned letters, five of five, and the controls read the same at eighty-seven of eighty-seven. The orders were reproduced at a second engine, equal at every type.
\chcert{1691}
\end{chproof}
\section{The cap reachable by genre}\label{mod:I-2-D534900-w-pgnc}
\begin{chwall}[The cap reached and left, type by type]\label{box:I-2-D534900-w-pgnc:1}
At the adopted passage reading every one of the six link types reaches its cap within the budget of the saturation, and a capped link leaves its cap under an admitted move at five of them and never at (4,4), so the types read open at (2,2), (2,3), (2,4), (3,3) and (3,4) and locked at (4,4); the readability order expects open at (2,2) and (2,3) and locked at the four others, and the two readings agree at three types only, parting first at (2,4). The open result stands at that parting: the dichotomy the saturation draws at this budget runs along the cut of a hub unit, not along the readability order.
\end{chwall}
\chstanding{open result}
\chdeps{Definitions 1.1 to 1.6 and Theorems 9.1 to 9.7 and Open result 9.8}{no theorem beyond it at this release}
The five adopted saturations are re-run under the admitted moves at the wave bound and the candidate bound, the cap of every type is read at the passage reading and at the valence reading beside, and two questions are put to the webs type by type, whether the cap is reached within the budget and whether a capped link leaves it under a move; the letters open and locked fall from those cells and are compared with the letters the readability order expects, the first parting type named with both cells; within the budget the cap is reached at the seeds themselves and held at the closed hubs of (4,4) alone, the leavings at the order-two types falling under the adopted cut.

\begin{chbridge}
The entanglement of a bipartite pure state is bounded by its Schmidt rank, at most the smaller of the two local dimensions, and the bound is reached at the maximally entangled states; whether a local operation keeps a state at the bound or lowers its rank is decided by the majorization order of its Schmidt coefficients \cite{Schmidt1907,Nielsen1999}. The link is the maximally entangled pair of two qudits and the tag is exact.
\end{chbridge}
\begin{chproof}[exact]
The premises are the link with its six types (chapter 1-2, a theorem), the moves in one table (chapter 1-4, a theorem) and the saturation with its seeds, waves and sinks (chapter 1-5, a theorem). The chain back to the axiom runs through chapter 1-2 and the framework chapter. From the six types, by the count of the records at both ends of a link of that type, we obtain the cap of every type,
\begin{equation}
c(2,2)=96,\quad c(2,3)=319,\quad c(2,4)=906,\quad c(3,3)=1036,\quad c(3,4)=2053,\quad c(4,4)=5100 .
\end{equation}
From the five saturations re-run under the twenty-eight admitted moves at the wave bound two and the candidate bound fifteen hundred per band, by the count of the links of each type that stand at their cap within the budget, we obtain the reaching cells, the capped links over the links of the type,
\begin{equation}
\begin{array}{lrrrrrr}
\text{type} & (2,2) & (2,3) & (2,4) & (3,3) & (3,4) & (4,4)\\
\text{links at the cap} & 92204 & 7530 & 197303 & 112591 & 71270 & 17146\\
\text{links of the type} & 94641 & 40530 & 259633 & 115027 & 72548 & 17656
\end{array}
\end{equation}
so every type is reached, the first reach at a seed at wave zero. From the same saturations, by the count of the capped links kept at a child with a passage removed at one of their ends (a leaving. A capped link removed whole by a move counted apart and never a leaving), we obtain the leaving cells,
\begin{equation}
\ell(2,2)=935,\quad \ell(2,3)=228,\quad \ell(2,4)=670,\quad \ell(3,3)=83,\quad \ell(3,4)=7,\quad \ell(4,4)=0 .
\end{equation}
By the letter rule, a type open when no link of the type reaches its cap within the budget or some capped link leaves it and locked when some link reaches and none leaves, we obtain the letters open at (2,2), (2,3), (2,4), (3,3) and (3,4) and locked at (4,4). Against the letters the readability order of chapter 1-2 expects, open at the order-one types (2,2) and (2,3) and locked at the four order-two types, the letters agree at (2,2), (2,3) and (4,4) and part at (2,4), (3,3) and (3,4). The first parting type is (2,4), its letter open against the expected locked, with the cap reached and left, which is the open result. The smallest instance on the page is the (2,2) link with its cap of ninety-six records, its first leaving is a cut of a capped link at a saturated site at the second wave of the bound-respecting seed, the link cut and the hub unit attached, so that the kept child carries the link below its cap.

The counter-tests alter the passage reading, at the valence reading beside, an end passed when the valence of its site is filled, the leaving counts become 935, 271, 1509, 1209, 789 and 110, every type reads open, the agreement with the readability order falls to two types and the first parting type is (2,4) again. The counter-tests of the derivation each moved the cell they were pinned to move and every control read the same. The reaching and leaving counts were reproduced at a second engine, equal at every cell.
\end{chproof}
\begin{figure}[h]\centering
\begin{adjustbox}{max width=\linewidth}\begin{tabular}{lrrr}\toprule
type of the link & near end & far end & cap \\ \midrule
(2,2) & 2 & 2 & 96 \\
(2,3) & 2 & 3 & 319 \\
(2,4) & 2 & 4 & 906 \\
(3,3) & 3 & 3 & 1036 \\
(3,4) & 3 & 4 & 2053 \\
(4,4) & 4 & 4 & 5100 \\ \bottomrule
\end{tabular}\end{adjustbox}
\caption{The link at its six types: the two end dimensions and the cap, the bound on the count of records a link of that type carries.}
\end{figure}
\section{The readout unit as the scheme at the one link}\label{mod:I-2-D535412-olus-t1}
\begin{chtheorem}[The readout unit as the scheme at the one link]\label{box:I-2-D535412-olus-t1:1}
At the one link the on-shell condition of the literature, residue one at the physical pole and the vertex at zero momentum transfer, is the unit lemma of the compendium, met as an exact identity by the construction of the readout unit and of the read map: the readout unit of the one link is an exact element at every derived cell, the trace form with the centre valued one, the identity reading one at the register dimension, the order step one tick at the pair state, the bare two-point function and the bare coupling both one; the read map product of the one-link record web is the identity on its labels with one pole at the unit, its residue the derived pole residue; and the residue is normalised at the readout unit at both registers with the vertex at zero displacement, the dressed residue being the coupling the scheme defines and never the normalisation it fixes.
\end{chtheorem}
\chstanding{theorem}
\chdeps{Definitions 1.1 to 1.6 and Theorems 9.1 to 9.7 and Open result 9.8}{no theorem beyond it at this release}
The readout unit of the one link is re-met as an exact element at every derived cell, the read map product of the one-link record web is shown the identity on its labels with the one pole at the unit and its residue equal to the derived pole residue, and the residue is normalized at the readout unit at both registers with the vertex at zero displacement; the scheme half of the third premise of the alpha correspondence thereby stands as a theorem at the one link, its bridging lemma drafted as text so that a Lean release may replace the opaque binder by the conjunction of the Ward statement and the on-shell scheme definition with its proof.

\begin{chbridge}
The on-shell renormalization condition of quantum electrodynamics fixes the normalization of the electron propagator at residue one at the physical pole and the vertex at zero momentum transfer, so that the renormalized charge is the charge measured at zero transfer; the Ward identity ties the vertex and the propagator renormalizations together \cite{Ward1950,BjorkenDrell1965}. The derivation supplies the two-point function of the one link at zero displacement, the readout unit and the pole with its residue here, and the scheme, the coupling, the momentum transfer and the Ward identity are named by correspondence.
\end{chbridge}
\begin{chproof}[finite and exact]
The premises are the pair ket and the transport to the far end (chapter 1-1, a theorem), the gain lemma with the Born count (chapter 1-3, a theorem) and the pole-residue result (chapter 5-1, a theorem). The chain back to the axiom runs through chapter 1-3 and the framework chapter. The procedure. The input is the one (2,3) link with its record web, one hundred seventy-two records at three hundred nineteen labels, and the readout unit, the trace form with the center valued one. For every label the read map is composed and its product over the web compared with the identity, the pole of the two-point function located and its residue read. The test applied is whether the product is the identity on the label and whether the residue, normalized at the readout unit, is one. The aggregation counts the labels at the identity, the pole and the residue at both registers. The exact part. From the readout unit, by the identity reading one at the register dimension and by the order step at the pair ket, we obtain the unit cells
\begin{equation}
\mathrm{tr}_{\text{unit}}(\mathbb{1}) = 1 \ \text{at } \dim R = 31104,\qquad \text{the order step } = 1 \text{ tick},\qquad \Pi_{\text{bare}} = 1,\qquad \alpha_{\text{bare}} = 1 .
\end{equation}
From the read map product over the one-link record web, by its evaluation on the three hundred nineteen labels, we obtain the identity with one pole at the unit,
\begin{equation}
\Pi(z) = \frac{319}{z-1},\qquad \operatorname{Res}_{z=1}\Pi = 319,\qquad \text{the two-read reading } 319^{2} = 101761 ,
\end{equation}
the residue equal to the derived residue of the pole-residue result at every reading, and from the residue over the readout unit at the label register and at the two-read register, we obtain the normalized residue one at both and the vertex at zero displacement with the momentum transfer zero, so the on-shell condition, residue one at the pole and the vertex at zero transfer, is met as an exact identity by the construction of the readout unit and of the read map, which is the theorem. The dressed residue, the coupling the scheme defines and never the normalization it fixes, stands beside as the exact element
\begin{equation}
Z_{3} = \frac{79 + 48\sqrt{2} + 42\sqrt{3} + 3\sqrt{6}}{31104} ,
\end{equation}
parting from one by construction. The smallest instance on the page is one label of the web, the read map composed around it reads one, so its product is the identity there and its pole residue over the unit is one.

The counter-tests struck the register dimension, the Killing form and the record count, each of which moved the cells it was pinned to move. read two reads in place of one, which moved no cell as pinned, and struck the text at its gate. Every control read the same at both engines, and the unit, the pole and the residue were reproduced at a second engine, equal at every cell.
\chcert{172}
\end{chproof}
\section{The soft limit at the one link}\label{mod:I-2-D535494-olsl-t1}
\begin{chconditional}[The soft limit at the one link]\label{box:I-2-D535494-olsl-t1:1}
Conditional on the six parent results of its union at their own standing, the pair ket fixed, the readout unit, the dressing equal to the bare, the two-point function of the one link, its pole residue and the momentum transfer at zero displacement.
The three clauses of the soft-limit structure are met at the one qutrit-qubit link, with no other link type, no hub, no composite body and no figure of the mass program in any cell: the one link is entered at the qutrit end with the far qubit end open at the flat readout and the unit Born square, re-met at the one-link cells and at every dock entry of that type; the washed factor at the open read and at zero displacement equals the readout unit from the pair state; and the two-point function of the record web with the far read washed equals the derived residue at the pole, unmoved. On this reading the second premise of the alpha correspondence is a theorem at the one link, with no clause parting.
\end{chconditional}
\chstanding{conditional theorem}
\chdeps{Definitions 1.1 to 1.6 and Theorems 9.1 to 9.7 and Open result 9.8}{no theorem beyond it at this release}
The wash (decoherence) of the open read at the far qubit end of the one link contributes the unit to the two-point function at the one pole, the same element at every reading of that far end, and the two-point function with the far read washed is the correlator of the record web multiplied by that unit with its residue the derived residue of the pole-residue result; so the soft-limit structure, one open read at the far end contributing a universal factor fixed at zero displacement, stands at the one link, conditional on the six parent results of its union at their own grades.

\begin{chbridge}
The soft photon theorem of Low and Weinberg: the amplitude for emitting a photon of vanishing momentum factorises into the amplitude without the photon multiplied by a universal factor fixed by the charges and momenta of the external lines, independent of the interior of the process \cite{Low1958,Weinberg1965}. The derivation supplies the washed read at the far end and the two-point function of the one link at zero displacement here, and the soft limit, the coupling fixed by it and the soft photon theorem are named by correspondence.
\end{chbridge}
\begin{chproof}[census]
The premises are the pair ket and the transport to the far end (chapter 1-1, a theorem), the readout unit at the one link (chapter 1-2, a theorem), the gain lemma with the Born count (chapter 1-3, a theorem), the measurement as the wash at one act (chapter 1-6, a theorem) and the pole-residue result (chapter 5-1, a theorem). The chain back to the axiom runs through chapter 1-2 and the framework chapter. By exhaustive enumeration of the one (2,3) link entered at its qutrit end with the far qubit end open, together with the fifty-three dock entries of that type at which the same far end is read (the set generated by the dock table of the saturation at the (2,3) type and by the one-link cells of the measurement act). The property tested being the three clauses of the soft-limit structure, factorization, universality and the limit at zero displacement. The counts by kind, the factorization clause met, the washed factor one against the Born square of the far marginal one, so the dressed two-point function with the far read washed is the two-point function of the one link multiplied by that factor,
\begin{equation}
\Pi_{\text{dressed}}(z) = \Pi(z)\cdot 1,\qquad \Pi(z) = \frac{319}{z-1} ;
\end{equation}
the universality clause met, the far end open at the (2,3) type with the readout one half at each outcome equal to the readout of the matched state of the measurement act, the factor a scalar of the far marginal alone, at fifty-three of fifty-three dock entries, and the limit clause met, the factor at zero displacement one against the factor at the open read one and the unit one, the washed residue three hundred nineteen against the derived residue three hundred nineteen at the pole one unmoved, and the momentum displacement zero. Three clauses met of three and no clause parting, so the second premise of the alpha correspondence is a theorem at the one link, which is the theorem, with no other link type, no hub, no composite body and no figure of the mass program entering any cell. The smallest case on the page is the one link itself, the far qubit read open at the flat readout gives the Born square one, the wash of that read leaves the unit, and the residue at the pole is unmoved.

The counter-tests flipped the open end, which moved five cells beginning at the open reading. left the read unwashed, which moved twelve cells beginning at the factor. reversed the link, which moved six cells. read the dock in place of the open end, which moved no cell as pinned, and struck the text at its gate. Every control read the same at both engines, the two site dimensions of every one of the one thousand six hundred ninety-one links of the web table read as an unordered pair, and the three clauses were reproduced at a second engine, equal at every cell.
\chcert{53}
\end{chproof}
\section{The record phase group at every genre}\label{mod:I-2-D535680-rphg-t1}
\begin{chtheorem}[The record phase group at every genre]\label{box:I-2-D535680-rphg-t1:1}
The record phase group is finite cyclic at every link type and every end: the phases of the admitted matchings, the entries and eigenvalues of the admitted cap bases and the stabiliser of the unit are roots of unity of the ring of the twenty-fourth root at every end, generating cyclic groups of order four at the qubit end, twelve at the qutrit end and eight at the ququart end, the (2,3) link at order twelve; every derived holonomy, the Berry result, the classical-limit sequences and the Chern result, lies inside its group; and the per-vector phases act by conjugation on every holonomy, so that the characteristic polynomial is free of every phase variable and the phase entry of the definition census is typed a gauge. The reading on which an antilinear holonomy is read as a linear matrix is printed beside it.
\end{chtheorem}
\chstanding{theorem}
\chdeps{Definitions 1.1 to 1.6 and Theorems 9.1 to 9.7 and Open result 9.8}{no theorem beyond it at this release}
The phases of every admitted matching are listed as monomials of the identification group at every end, the admitted cap bases are read at their entries and eigenvalues with the certified Lean bases re-met, and every phase is shown a root of unity generating a finite cyclic group end by end and type by type; every derived holonomy result is recomputed at the adopted transport and placed inside its group; the conjugation lemma is proved cycle by cycle at the adopted reading.

\begin{chbridge}
The geometric phase of Berry and the phase of Pancharatnam that precedes it: a state transported around a closed loop of parameters returns to itself up to a phase fixed by the loop alone, and for a cyclic sequence of projective measurements the phase is the argument of the product of the overlaps around the cycle \cite{Berry1984,Pancharatnam1956}. The derivation supplies the phase group and the holonomy of every cycle here, and the Berry phase and the Chern number are named by correspondence.
\end{chbridge}
\begin{chproof}[census]
The premises are the pair ket and the transport to the far end (chapter 1-1, a theorem), the link with its six types (chapter 1-2, a theorem) and the orbit reading of the knowledge clause with its identification group (chapter 1-1, a conditional theorem). The chain back to the axiom runs through chapter 1-1 and the framework chapter. By exhaustive enumeration of the phases of the admitted matchings (the set generated by the fourth-root monomials of the identification group at its orders thirty-two, three hundred eighty-four and six thousand one hundred forty-four at the three end dimensions), of the entries and eigenvalues of the twelve admitted cap bases, with the four certified bases of the package re-met, and of the stabilizer of the unit. The property tested being whether every phase is a root of unity in the ring of the twenty-fourth root and which cyclic group the phases of an end generate. The counts by kind, every phase a root of unity with no non-root element, the groups generated cyclic of order four at a qubit end, twelve at a qutrit end and eight at a ququart end, and by link type
\begin{equation}
|G_{(2,2)}| = 4,\quad |G_{(2,3)}| = 12,\quad |G_{(2,4)}| = 8,\quad |G_{(3,3)}| = 12,\quad |G_{(3,4)}| = 24,\quad |G_{(4,4)}| = 8 ,
\end{equation}
the (2,3) link at order twelve, the fourth roots of its qubit end joined with the cube roots of its qutrit end, and the twenty-fourth roots reached at the join of the matching phases with the ray-cycle phases of the cap bases. Every derived holonomy inside its group, the Berry result recomputed from its dimension sequences at the transport taken equal to the derived spectra and inside at one thousand one hundred eighty-three of one thousand one hundred eighty-three, the classical-limit sequences equal at twenty-four of twenty-four with their closed paths at one thousand one hundred forty-four of one thousand one hundred forty-four and their walks inside at two hundred ten of two hundred ten, and the Chern result at nine hundred seventy-two defined spectra inside at every one. The conjugation lemma proved at one hundred two of one hundred two cycles, the per-vector phases a diagonal unitary at every end acting by conjugation on every linear holonomy and on the linear element every antilinear holonomy defines, the characteristic polynomial free of every phase variable. The first parting element none, so the group is finite cyclic at every type and the phase entry of the definition census is a gauge, which is the theorem. The smallest case on the page is a qubit end, its matching phases are the fourth roots of unity,
\begin{equation}
\{1,\ i,\ -1,\ -i\} \cong \mathbb{Z}/4 ,
\end{equation}
and the holonomy of the (2,2) link lies inside that group.

The counter-tests admitted the eighth roots alone, which moved nine cells beginning at the generated groups. offered a continuous phase variable at an entry, which the derivation refused as no root of unity of the ring, as pinned. made one end non-diagonal, which moved no cell. struck the conjugate, which moved eleven cells beginning at the lemma, and struck the text at its gate. Every control read the same at both engines, and the groups and the holonomies were reproduced at a second engine, equal at every cell. At the reading beside, an antilinear holonomy read as a linear matrix parts the lemma at twenty-three of the one hundred two cycles, printed beside for comparison.
\chcert{1183}
\end{chproof}
\section{The cap under the derived definition of open}\label{mod:I-2-D539697-w-aptc}
The entanglement a pure bipartite state carries is bounded by the smaller local dimension: the Schmidt rank of a pair of a qudit of dimension d with a qudit of dimension d prime is at most the smaller of the two, so a pair of unequal dimensions is maximally entangled only on a subspace of the larger factor, and the remaining dimensions of that factor are reached by no term of the state \cite{Nielsen2000,Horodecki2009}. The link and its six genres are introduced at section~\ref{mod:I-2-link-and-its-six-genres} and the saturation with its sinks at section~\ref{mod:I-5-saturation-in-one-figure-seeds-waves-termination-sinks}.
Under the derived definition of open (section~\ref{mod:I-1-D539391-w-dole}) every link is typed by its end dimensions alone: an unequal link is uncapped at its far end set and an equal link is capped at both ends. Read at the cap of every body, this typing parts from the body-alone typing at the first of its clauses, the reach bound.
The body-alone typing reads every lepton link uncapped by its genre letters and every nucleon link capped by the cap census, while the derived typing reads the (3,3) core links and the (2,2) pairs capped and the (2,3) docks and the hub legs uncapped at every body alike. Against the register of the reach result, the reach bound under the derived typing agrees at none of the five bodies, the first parting at the electron, 2078 against 79; the strike at the uncapped entry agrees at two of the five bodies, the first parting at the tau, 0 against 8; the dock link of the electron ladder reads uncapped under the derived definition at the register 175 against capped at the register 319; and the cap census of the universe web under the derived definition terminates at wave 0 with 221 links capped and 330 uncapped, against the cap wave of the standing census at wave 1 with 497 capped, the derived partition agreeing with the depth census at the link level at 339 of 339 listed links. The bound-respecting result at the proton and the sink-weight result at the electron and the muon are reproduced, and two of the six clauses are re-derived.
The result is an open one, an open result: the stronger statement, that the derived letters agree with the body-alone typing at every link of every body, is not met, and the result stands at the weaker reading with its first clause, the reach, named.
What stays open is a typing under which the derived letters agree with the body-alone typing at every link of every body, found at none; the reach result stands beside at its register corollary.
\section{The six genres against the terms of the Standard Model}\label{mod:I-2-D540373-the-six-genres-against-the-lagrangian-terms}
\begin{chtable}[The six genres against the terms of the Standard Model]\label{box:I-2-D540373-the-six-genres-against-the-lagrangian-terms:1}
This instrument keys the six genres, the six unordered pairs of the three registers, to their register pair, their first non-vanishing read order, their carrier move classes and their identification, six of six complete, and lays one sentence at each of the twelve slots of the pair ket and the attachment citing its object by its row, its table and its cell, twelve of twelve, with every figure a cited cell and every cited cell located at one hundred four of one hundred four; the term of the Standard Model stands at a column in the vocabulary of the literature, and the two slots beyond and the three points at which the correspondence is not one to one, the hypercharge a relation of two genres, the self-term and the coupling at one genre and the theta term at no row, are printed as sentences.
\end{chtable}
\chstanding{instrument}
\chdeps{Definitions 1.1 to 1.6, Theorems 9.1 to 9.7 and Open result 9.8 and Open results 241.1 and 241.2, Theorems 241.3 and 241.4 and Open results 241.5 to 241.7}{no theorem beyond it at this release}
The six genres are keyed from the summary of the register pairs, the read order of the readability theorem, the carrier move classes of the forces open result and the identification of the census; each slot sentence is checked at its cited cells, and the instrument is complete when every slot cites its object and carries no figure beyond its cells.

\begin{chbridge}
The Lagrangian of the Standard Model is a sum of terms, the gauge kinetic terms of the three interactions, the fermion kinetic terms with their covariant derivatives, the Yukawa couplings, the Higgs potential with its self-coupling and the theta term \cite{Glashow1961,Weinberg1967}. The derivation supplies the genres with their read orders and move classes here, and the terms are named by correspondence.
\end{chbridge}
\begin{chproof}[instrument]
The table has one row per genre, six rows, and the columns the register pair, the first non-vanishing read order at the adopted face of the readability theorem, the carrier move classes present at the move graph of the sinks, the identification of the solution reading and the names the readability row attaches. The rule that produced it: a genre is keyed when every column carries its cell by locator from the parent table, and a slot is complete when its sentence cites its object by row, table and cell with no figure beyond the cited cells.
\begin{equation}
\begin{array}{llll}
\text{genre} & \text{read order} & \text{the carrier move classes} & \text{identification}\\
(2,2) & 1 & \text{hub-forming, join, cut, hub step} & \text{electromagnetic}\\
(2,3) & 1 & \text{hub-forming, cut, hub step} & \text{electromagnetic}\\
(2,4) & 2 & \text{hub-forming, cut, hub step} & \text{weak}\\
(3,3) & 2 & \text{none} & \text{strong}\\
(3,4) & 2 & \text{hub-forming, cut, hub step} & \text{the Higgs}\\
(4,4) & 2 & \text{cut, hub step} & \text{the Higgs}
\end{array}
\end{equation}
The twelve slots beneath, the pair ket and the attachment at each genre, each carry one sentence citing their object at its row and cell, twelve of twelve complete, the first incomplete none. The three points at which the correspondence is not one to one are printed as sentences: the hypercharge term is a relation of the electromagnetic identification at the (2,2) genre and the weak identification at the (2,4) genre and stands at no one slot; the self-term and the coupling of the Higgs fall at the one (4,4) genre, the quartic the Higgs office over the vacuum office squared and the couplings ratios of the same office at the same vacuum; and the theta term stands at no row of the table, the strong genre carried by no move class present at the sinks.

The counter-tests struck the office table of the photon row, which left the (2,2) pair-ket slot unkeyed and the instrument incomplete with that slot first. A second refused a wrong hash at the gate. A third struck the item at the unread slots, the (3,4) pair ket first, and a fourth struck the term column at seventy-nine terms, which left every status unchanged. Every control read the same at both engines, eight of eight, and the table was re-cut twice, equal.
\end{chproof}
\section{The genre functions in the end dimensions}\label{mod:I-2-D542526-gnfn-t1}
\begin{chtheorem}[The genre functions in the end dimensions]\label{box:I-2-D542526-gnfn-t1:1}
At the six genres, (2,2), (2,3), (2,4), (3,3), (3,4) and (4,4), with $f$ the smaller end and $c$ the larger, every structural cell of the matrix is a closed form in the end dimensions, met at six of six genres and six of six columns: the identification algebra has the dimension $f^{2} + c^{2} - 1$, seven, twelve, nineteen, seventeen, twenty-four and thirty-one; the central scalar is $c - f$, zero, one, two, zero, one and zero; the register of the count is the cut multiplied by the whole multiplied by the whole, $2304$, $31104$, $184320$, $419904$, $2488320$ and $14745600$; the first non-vanishing read order is one at an end product at most six and two beyond it, one, one, two, two, two and two; the office of the pair ket under the body-alone rule is zero at the struck dock and pair genres and the act squared multiplied by one plus the squared read at order two at the dark genres; and the stabiliser of the pair ket has the dimension $f^{2}$ plus the special unitary dimension of the difference, four, four, seven, nine, nine and sixteen. The first column with no closed form is none.
\end{chtheorem}
\chstanding{theorem}
\chdeps{Definitions 1.1 to 1.6, Theorems 9.1 to 9.7 and Open result 9.8, Theorem 124.1, Theorems 185.1 to 185.12 and Theorems 200.1 to 200.3 and Open result 200.4}{no theorem beyond it at this release}
The matrix is built from the pair ket at the end dimensions of each genre, the identification algebra as an explicit basis of block-diagonal anti-Hermitian traceless matrices with its brackets closed and its centre one line, the central scalar on the pair ket, the register of the count from the whole and cut registers, the read order under the count law, the office under the body-alone rule and the stabiliser as the null space of the tensor action at the pair ket; the closed forms, pre-declared ahead of the run, are evaluated at every genre and compared cell for cell.
\begin{figure}[htbp]\centering
\scriptsize
\setlength{\tabcolsep}{2pt}
\begin{adjustbox}{max width=\linewidth}\begin{tabular}{|p{13mm}|p{11mm}|p{9mm}|p{12mm}|p{11mm}|p{33mm}|p{12mm}|p{19mm}|p{22mm}|p{13mm}|}
\hline
\textbf{genre, the smaller end $f$ and the larger $c$} & the identification algebra dimension \newline $f^{2}+c^{2}-1$ & the central scalar \newline $c-f$ & the register of the count \newline $cut\cdot whole\cdot whole$ & the first non-vanishing read order \newline $1+\mathrm{ind}(fc>6)$ & the office of the pair ket under the body-alone rule \newline $\mathrm{ind}(fc>6)\cdot (f\cdot c)^2\cdot (1+G_{2}^{2})$ & the stabiliser dimension at the pair ket \newline $f^{2}+(c-f)^{2}-\mathrm{ind}(c>f)$ & the bound at the pair ket, the family maximum at distinct bases and free labels & the bound at the pair ket, the union maximum over every named sub-family & the projected stabiliser at the larger end \\ \hline
\textbf{(2,2)}, the index \newline $f=2$, $c=2$ & $7$ & $0$ & $2304$ & $1$ & $0$ & $4$ & 2 & 2 root 2 & $4$ \\ \hline
\textbf{(2,3)} \newline $f=2$, $c=3$ & $12$ & $1$ & $31104$ & $1$ & $0$ & $4$ & (3 plus root 3) over 3 & 2 root 2 & $4$ \\ \hline
\textbf{(2,4)} \newline $f=2$, $c=4$ & $19$ & $2$ & $184320$ & $2$ & $466566137/7290000$, $2/6075$, $0$, $0$ & $7$ & 2 & 2 root 2 & $7$ \\ \hline
\textbf{(3,3)} \newline $f=3$, $c=3$ & $17$ & $0$ & $419904$ & $2$ & $6879719125/84934656$, $0$, $1675/21233664$, $0$ & $9$ & 2 & (12 plus 8 root 3) over 9 & $9$ \\ \hline
\textbf{(3,4)} \newline $f=3$, $c=4$ & $24$ & $1$ & $2488320$ & $2$ & $4299817692073/\allowbreak 29859840000$, $0$, $5069/\allowbreak 373248000$, $0$ & $9$ & (2 plus root 3) over 2 & (12 plus 8 root 3) over 9 & $9$ \\ \hline
\textbf{(4,4)} \newline $f=4$, $c=4$ & $31$ & $0$ & $14745600$ & $2$ & $46656001681/182250000$ & $16$ & 2 & (2 root 2 plus 6 c minus 2 root 2 c) over 3 & $16$ \\ \hline
\end{tabular}\end{adjustbox}
\caption{The genre-functions matrix in the end dimensions, with $f$ the smaller end and $c$ the larger: every structural cell of the matrix of the six genres is a closed form in the end dimensions, the identification algebra dimension, the central scalar, the register of the count, the first non-vanishing read order, the office of the pair ket under the body-alone rule and the stabiliser dimension at the pair ket, met at six of six genres and six of six columns; beside the six columns, the bound at the pair ket of each genre at its two readings, the family maximum at distinct bases and free labels and the union maximum over every named sub-family (the Bell bound of the literature, the Weyl-family maximum of the first Bell open result and the every-basis union maximum of the second), and the projected stabiliser at the larger end, the stabiliser of the pair ket of the link projected to the larger end; the (2,2) genre is marked as the index. A cell of several parts is an element of the ring in the basis one, root two, root three and root six, its parts in that order.}\label{fig:I-2-D542526-genre-functions-matrix}
\end{figure}

\begin{chbridge}
A state of two systems of dimensions $f$ and $c$, with $f$ at most $c$, has a Schmidt decomposition with at most $f$ terms, and a maximally entangled state is fixed by the unitaries of the larger system that act as the transpose of a unitary of the smaller on its support and trivially on the complement \cite{Schmidt1907,Horodecki2009}. The derivation supplies the six columns in the end dimensions here, and the genres carry the names of the literature by correspondence.
\end{chbridge}
\begin{chproof}[finite and exact]
The premises are the readability order at the six genres (chapter 1-2, a theorem), the count law at the six genres (chapter 8-1, a theorem), the Higgs and carrier offices with the body-alone rule (chapter 6-10, a theorem) and the stabiliser of the pair ket (chapter 5-1, an open result). The chain back to the axiom runs through chapter 1-2 and the framework chapter. The procedure. The input is the pair ket at each of the six genres, the maximally entangled vector of the two ends of dimensions $f$ and $c$. The enumeration rule builds, at each genre, the identification algebra as the block-diagonal traceless anti-Hermitian matrices on the direct sum of the two ends, tests the closure of every bracket, reads the centre, acts on the pair ket by the tensor action and reads the null space; the aggregation prints the six columns and evaluates the six pre-declared forms. From the two ends, by the block construction, we obtain the algebra of dimension $f^{2} - 1 + c^{2} - 1 + 1 = f^{2} + c^{2} - 1$, the brackets closed at twenty-one, sixty-six, one hundred seventy-one, one hundred thirty-six, two hundred seventy-six and four hundred sixty-five pairs, the centre one line with the entry $c$ at the $f$ end and $-f$ at the $c$ end. From the central generator on the pair ket of the matched directions, by the action, we obtain the central scalar $c - f$. From the whole register $f^{2} c^{2}$ and the cut register $(f^{2} - 1)(c^{2} - 1)$, by the register of the count at order one, the cut multiplied by the whole multiplied by the whole,
\begin{equation}
R = (f^{2} - 1)(c^{2} - 1)\, f^{4} c^{4} ,
\end{equation}
two thousand three hundred four at (2,2) and fourteen million seven hundred forty-five thousand six hundred at (4,4). From the sectors of the count law, by the complete-read bound six, we obtain the read order one if and only if the end product is at most six and two otherwise, and the office of the pair ket zero at the struck (2,2) and (2,3) links and, at the dark genres, the act squared multiplied by one plus the squared read at order two,
\begin{equation}
\omega_{(2,4)} = \tfrac{466566137}{7290000} + \tfrac{2}{6075}\sqrt{2} , \qquad \omega_{(4,4)} = \tfrac{46656001681}{182250000} ,
\end{equation}
with the reads $\tfrac{1}{288} + \tfrac{1}{1350}\sqrt{2}$ at (2,4), $\tfrac{67}{82944} + \tfrac{25}{41472}\sqrt{3}$ at (3,3), $\tfrac{137}{414720} + \tfrac{37}{259200}\sqrt{3}$ at (3,4) and $\tfrac{41}{216000}$ at (4,4). From the tensor action at the pair ket, by the null space, we obtain the stabiliser of dimension $f^{2} + (c - f)^{2}$, less one where $c$ exceeds $f$, the ranks three, eight, twelve, eight, fifteen and fifteen. Every form evaluated at the six genres equals the computed cell at thirty-six of thirty-six cells, the first miss none, which is the theorem. The smallest instance on the page is the (2,2) genre: the seven of the algebra with the central generator
\begin{equation}
\mathrm{diag}(2, 2, -2, -2) ,
\end{equation}
the central scalar zero, the register two thousand three hundred four, the read order one, the office zero at the struck pair link and the stabiliser of dimension four.

The counter-tests inserted a pair with an end of dimension five, which moved every column with the forms still equal to the computed cells at twenty-four of twenty-four. A second exchanged the end dimensions, which moved the sign of the central scalar at the unequal genres. A third replaced the pair ket by a product ket, which moved the stabiliser. A fourth struck the centre, which moved the algebra dimension by one and struck the scalar, and a fifth struck the text at its gate, which left the gate unmet. Every control read the same at both engines, twelve of twelve, and the matrix was reproduced at a second engine, equal at every cell.
\chcert{36}
\end{chproof}
\begin{figure}[h]\centering
\begin{adjustbox}{max width=\linewidth}\begin{tabular}{lrrrrr}\toprule
genre & algebra & central scalar & register & read order & stabiliser \\ \midrule
(2,2) & 7 & 0 & 2304 & 1 & 4 \\
(2,3) & 12 & 1 & 31104 & 1 & 4 \\
(2,4) & 19 & 2 & 184320 & 2 & 7 \\
(3,3) & 17 & 0 & 419904 & 2 & 9 \\
(3,4) & 24 & 1 & 2488320 & 2 & 9 \\
(4,4) & 31 & 0 & 14745600 & 2 & 16 \\ \bottomrule
\end{tabular}\end{adjustbox}
\caption{The matrix of the six genres at five of its six columns, each a closed form in the end dimensions $f$ and $c$: the identification algebra $f^{2} + c^{2} - 1$, the central scalar $c - f$, the register of the count, the read order and the stabiliser of the pair ket.}
\end{figure}
\section{The tensor type at every genre}\label{mod:I-2-D543154-w-tntp}
\begin{chwall}[The tensor type at every genre]\label{box:I-2-D543154-w-tntp:1}
At the six genres of the universe web the local increment operator of the serial insertion at the washed Laplacian is the same local one-link difference at every genre, copy number and rung, the cells zero, minus one, one; minus one, two, minus one; one, minus one, zero on the three support sites in the chain order, the trace two, the rank two, the spectrum minus one, zero and three, symmetric and even under the end swap, with two trivial and two standard pieces under the site permutations, a scalar of the rotation group at the site index under the adopted lift at every genre, the five-dimensional component zero at one hundred twenty-six of one hundred twenty-six cells, while the wash form factor of the curvature row is one form for every genre whose fixed forms carry a nonzero five-dimensional component, so every one of the six genres carries the symmetric traceless rank-two type at its wash form alike. The first face, the (2,2) genre the one whose increment or wash form carries the type and the other five not, reads no, with the first parting the (2,3) genre at the wash form, the other genres carrying at five of five; the second face, another genre carrying it, is named at its first branch. At the far face the type is present at the (2,3) pair by the directions table and at the (2,2) pair under the transport reading alone, and at no other genre, no row defining a far face there.
\end{chwall}
\chstanding{open result}
\chdeps{Definitions 1.1 to 1.6, Theorems 9.1 to 9.7 and Open result 9.8, Theorem 53.1, Theorems 90.1 to 90.4 and Theorem 115.1}{no theorem beyond it at this release}
The chain census of the universe web is taken at every genre under the chain rule of the serial row and the serial insertion is applied at one to three copies and seven rungs at each; the increment, the difference of the graph Laplacians at a step, is computed on its three-site support, decomposed under the site permutations and under the lifted rotation group at the adopted lift, and the wash form factor of the curvature row and the far-end object of every genre are decomposed alike.

\begin{chbridge}
The perturbation of a metric around a background decomposes under rotations into a scalar, a vector and a symmetric traceless tensor \cite{Weinberg1965,Feynman1963}. The derivation supplies the increment of every genre and its decomposition here. The correspondence is by name.
\end{chbridge}
\begin{chproof}[finite and exact]
The premises are the serial (2,2) increment at the washed Laplacian (chapter 10-6, an open result), the rotation group at the increment (chapter 10-0, an open result), the curvature as the second difference (chapter 10-0, an open result), the rotation group (chapter 3-2, a theorem) and the Lorentz group at the pair (chapter 10-0, a theorem). The chain back to the axiom runs through chapter 1-2 and the framework chapter. The procedure. The input is the universe web of the hand-off, four hundred forty-four sites and five hundred fifty-one links at the genres (2,2) eighty-five, (2,3) thirty-four, (2,4) two hundred twenty-nine, (3,3) one hundred thirty-four, (3,4) sixty-seven and (4,4) two. The enumeration rule takes the chain census at every genre under the chain rule of the serial row, the letters by the component of the cut, with the adopted chain the first bulk link in the link order, applies the serial insertion at one to three copies at every genre, the new site between the previous site of the chain run and the last site at the dimension that makes the new link a link of the genre, and computes the increment as the difference of the graph Laplacians at every copy number and at rungs zero to six; the test at each cell is equality of the increment with the one-link difference; the aggregation counts the cells. From the insertion at any genre, copy number and rung, by the difference of the Laplacians, we obtain the increment on the three support sites in the chain order,
\begin{equation}
\Delta = \begin{pmatrix} 0 & -1 & 1 \\ -1 & 2 & -1 \\ 1 & -1 & 0 \end{pmatrix} ,
\end{equation}
the trace two, the row sums zero, the rank two, the characteristic polynomial $x^{3} - 2x^{2} - 3x$ with the spectrum minus one, zero and three, the same at one hundred twenty-six of one hundred twenty-six cells, at the (2,2) genre re-met against the increment table of the serial row at one hundred five of one hundred five. From the increment, by the site permutations, we obtain the trace part with the entries two thirds on the diagonal and the traceless part
\begin{equation}
\Delta_{0} = \begin{pmatrix} -\tfrac{2}{3} & -1 & 1 \\ -1 & \tfrac{4}{3} & -1 \\ 1 & -1 & -\tfrac{2}{3} \end{pmatrix} ,
\end{equation}
of rank three, two trivial pieces and two standard pieces and no sign piece. From the increment extended by the identity at the factor of the record space, by the adjoint action of the lifted rotation generators at the adopted lift, the orientation sign one at every end, we obtain the action zero at three of three with the Casimir zero and the five-dimensional component zero at every genre, copy number and rung, a scalar at the site index. From the wash form factor of the curvature row, by the same decomposition under the four-by-four generators, we obtain its two fixed forms, the diagonal one, minus two, minus two, zero and the diagonal zero, zero, zero, one, each with a nonzero five-dimensional component, and the whole form of the trivial type alone, so that the washed increment of every genre carries the type at its wash form, at the first two forms yes and at the whole form no, the same at the six genres. The far-end object of every genre is a scalar at the ends reading, and at the far face of the (2,3) pair by the directions table and at the far end of the (2,2) pair under the transport reading it carries the five-dimensional component with the weights zero and two at the piece of dimension six, no far face being defined at the other four genres. So the (2,2) genre is one of six and not the one, the first other genre the (2,3) at the wash form, which is the open result. The smallest instance on the page is the one-link reference, the Laplacian of one link with the trace two, the rank one and the spectrum zero and two, from which the increment is the two new one-link Laplacians less the removed one.

The counter-tests struck the wash to the unit form on the four coordinates, which moved seventeen cells with no genre carrying at the wash form. A second put the (2,3) increment in place of the (2,2) as the distinguished genre, which moved three cells. A third lifted the rotation group at the site direction, which moved nine cells with every genre carrying at the increment. A fourth transposed the increment, which moved nothing, the increments symmetric, and a fifth struck the text at its gate, which left the gate unmet. Every control read the same at both engines, fourteen of fourteen, and the increments were reproduced at a second engine, equal at every cell.
\chcert{126}
\end{chproof}
\begin{figure}[h]\centering
\begin{adjustbox}{max width=\linewidth}\begin{tabular}{lrlll}\toprule
genre & links & the type at the increment & the type at the wash form & the type at the far face \\ \midrule
(2,2) & 85 & no & yes & yes \\
(2,3) & 34 & no & yes & yes \\
(2,4) & 229 & no & yes & no face defined \\
(3,3) & 134 & no & yes & no face defined \\
(3,4) & 67 & no & yes & no face defined \\
(4,4) & 2 & no & yes & no face defined \\ \bottomrule
\end{tabular}\end{adjustbox}
\caption{The carrying census of the symmetric traceless rank-two type at the six genres of the universe web: at the site index of the increment at none, at the wash form at every one, at the far face at the two pairs where a far face is defined.}
\end{figure}
\par\medskip\noindent{\itshape The genres fixed, the thread turns to what a read leaves behind along a link, the records and the count that weights them.}\label{narr:close:1-2}\par
\section{Records, reads and the Born count: the gain lemma}\label{mod:I-3-records-reads-and-the-born-count-the-gain-lemma}
\chdesig{1-3}{}{I-3-records-reads-and-the-born-count-the-gain-lemma}{I-3-D540846-phjr-t1}
\label{narr:open:1-3}A read leaves records, and the count of records at a sector is what the rest of the compendium weights with. Theorem~18.2, the gain lemma, fixes that weight at the spectrum of records; Theorem~18.1 gives the read-state boundary, Theorem~18.3 the gain as the washed direction count at a sector, and Theorem~18.4 the boson as the chain readout. The open results of the chapter mark the edges of the read: the soft limit as the washed open read, Open result~19.1, the dressing of what the read does not pass, Open result~20.1, the stop at the cap, Open result~21.1, the enclosed record count read as a phase, Open result~22.1, and the geodesic as the passage, Open result~23.1, with the three readout rules read as one read at three levels in Theorem~24.1.
\begin{chconditional}[The read-state boundary]\label{box:I-3-records-reads-and-the-born-count-the-gain-lemma:1}
Conditional on the results this theorem rests on at their own standing, as its derivation names them.
This chapter establishes that the Born count is a measure and the readout its root. Every clause of the six lines of the text types by its head as a read-side or a state-side clause, every state quantity the text forces is separated by the readouts of the admitted reads, and every quantity it names without forcing leaves every one-site readout unchanged under the wash (decoherence), so the text forces nothing about a state beyond what its reads fix.
\end{chconditional}
\chstanding{conditional theorem}
\chdeps{Definitions 1.1 to 1.6, Theorems 2.1 to 2.3, Open results 2.4 and 2.5, Theorem 2.6, Open results 2.7 and 2.8 and Theorems 2.9 and 2.10 and Theorems 9.1 to 9.7 and Open result 9.8}{the theorems of Chapters 1-1, 1-3, 1-5 to 1-6, 3-1 to 3-4, 3-6 to 3-8, 5-1, 5-3, 5-5 to 5-6, 5-8 to 5-9, 6-2, 6-4 to 6-5, 6-10, 7-1 to 7-2, 7-4 to 7-7, 8-3, 9-5, 10-1 to 10-2, 11-1, 12-2, 13-1, 14-2, 14-6, 15-1, 15-3 to 15-6}
\begin{chtheorem}[The gain lemma at the spectrum of records]\label{box:I-3-records-reads-and-the-born-count-the-gain-lemma:2}
The multiplicity of every record is the product of the dimensions of the ends it passes, one at no end, two at the qubit end, three at the qutrit end and six at both, at all one hundred seventy-two records of the qutrit-qubit link, its four end-set sums agreeing with the support-rank result; the record web therefore carries a link-product measure equal to the Born count at every end set and every link type.
\end{chtheorem}
\chstanding{theorem}
\chdeps{Definitions 1.1 to 1.6, Theorems 2.1 to 2.3, Open results 2.4 and 2.5, Theorem 2.6, Open results 2.7 and 2.8 and Theorems 2.9 and 2.10 and Theorems 9.1 to 9.7 and Open result 9.8}{the theorems of Chapters 1-1, 1-3, 1-5 to 1-6, 3-1 to 3-4, 3-6 to 3-8, 5-1, 5-3, 5-5 to 5-6, 5-8 to 5-9, 6-2, 6-4 to 6-5, 6-10, 7-1 to 7-2, 7-4 to 7-7, 8-3, 9-5, 10-1 to 10-2, 11-1, 12-2, 13-1, 14-2, 14-6, 15-1, 15-3 to 15-6}
\begin{chconditional}[Gain as the washed direction count at the sector]\label{box:I-3-records-reads-and-the-born-count-the-gain-lemma:3}
Conditional on the results this theorem rests on at their own standing, as its derivation names them.
The gain squared of a record is the product, over the signed exchanges a record survives, of the dimension of the end each exchange acts at; the gain is also the product of two tests at the two ends by the commutant of each end; and under the definition of open the three per-end factors are the three transport statuses, the near end one, the open far end the dimension less one and the washed far end the dimension plus one, so the role rule is a corollary of that definition and the one posed assumption of the fine-structure chain is discharged.
\end{chconditional}
\chstanding{conditional theorem}
\chdeps{Definitions 1.1 to 1.6, Theorems 2.1 to 2.3, Open results 2.4 and 2.5, Theorem 2.6, Open results 2.7 and 2.8 and Theorems 2.9 and 2.10 and Theorems 9.1 to 9.7 and Open result 9.8}{the theorems of Chapters 1-1, 1-3, 1-5 to 1-6, 3-1 to 3-4, 3-6 to 3-8, 5-1, 5-3, 5-5 to 5-6, 5-8 to 5-9, 6-2, 6-4 to 6-5, 6-10, 7-1 to 7-2, 7-4 to 7-7, 8-3, 9-5, 10-1 to 10-2, 11-1, 12-2, 13-1, 14-2, 14-6, 15-1, 15-3 to 15-6}
\begin{chconditional}[The boson as the chain readout]\label{box:I-3-records-reads-and-the-born-count-the-gain-lemma:4}
Conditional on the results this theorem rests on at their own standing, as its derivation names them.
The carrier webs read as chains give the bosons as chain readouts, a theorem conditional on the launch-weight rule and the closure construction of its own study.
\end{chconditional}
\chstanding{conditional theorem}
\chdeps{Definitions 1.1 to 1.6, Theorems 2.1 to 2.3, Open results 2.4 and 2.5, Theorem 2.6, Open results 2.7 and 2.8 and Theorems 2.9 and 2.10 and Theorems 9.1 to 9.7 and Open result 9.8}{the theorems of Chapters 1-1, 1-3, 1-5 to 1-6, 3-1 to 3-4, 3-6 to 3-8, 5-1, 5-3, 5-5 to 5-6, 5-8 to 5-9, 6-2, 6-4 to 6-5, 6-10, 7-1 to 7-2, 7-4 to 7-7, 8-3, 9-5, 10-1 to 10-2, 11-1, 12-2, 13-1, 14-2, 14-6, 15-1, 15-3 to 15-6}
\begin{chwall}[Role from the arrow passage]\label{box:I-3-records-reads-and-the-born-count-the-gain-lemma:5}
Nine open results stand beside. The role rule is not re-derived from the arrow passage; the orbit norm of a record word is not the gain, the orbit sum vanishing at every odd and mixed-even record; the inverse wash overlap is not the gain at any reading; the sector measure identity is arithmetic at every term and a clause at none; the Born count as the dressed measure at the full web parts from the derived figure at its deciding cell; the readout rule is not one output at every mass entry, parting first at the muon; the fine-structure condition is not discharged at the second read; and the readout rules at the link weights and at the sector weights re-meet no located entry, parting first at the electron.
\end{chwall}
\chstanding{open result}
\chdeps{Definitions 1.1 to 1.6, Theorems 2.1 to 2.3, Open results 2.4 and 2.5, Theorem 2.6, Open results 2.7 and 2.8 and Theorems 2.9 and 2.10 and Theorems 9.1 to 9.7 and Open result 9.8}{the theorems of Chapters 1-1, 1-3, 1-5 to 1-6, 3-1 to 3-4, 3-6 to 3-8, 5-1, 5-3, 5-5 to 5-6, 5-8 to 5-9, 6-2, 6-4 to 6-5, 6-10, 7-1 to 7-2, 7-4 to 7-7, 8-3, 9-5, 10-1 to 10-2, 11-1, 12-2, 13-1, 14-2, 14-6, 15-1, 15-3 to 15-6}
\begin{chwall}[The readout law along the sector route of records]\label{box:I-3-records-reads-and-the-born-count-the-gain-lemma:6}
Three further open results stand at this chapter, each at its grade: the readout law along the sector route of records, the order-one result at the web record spectrum, and the wash integer at the fan loops of the enclosed set.
\end{chwall}
\chstanding{open result}
\chdeps{Definitions 1.1 to 1.6, Theorems 2.1 to 2.3, Open results 2.4 and 2.5, Theorem 2.6, Open results 2.7 and 2.8 and Theorems 2.9 and 2.10 and Theorems 9.1 to 9.7 and Open result 9.8}{the theorems of Chapters 1-1, 1-3, 1-5 to 1-6, 3-1 to 3-4, 3-6 to 3-8, 5-1, 5-3, 5-5 to 5-6, 5-8 to 5-9, 6-2, 6-4 to 6-5, 6-10, 7-1 to 7-2, 7-4 to 7-7, 8-3, 9-5, 10-1 to 10-2, 11-1, 12-2, 13-1, 14-2, 14-6, 15-1, 15-3 to 15-6}
The six lines are typed clause by clause at their tokens and every derived readout construction is placed in the sector the reads fix. Every record of the one link is a link at the dimensions of its passed ends, so its multiplicity is their product, and the four end-set sums and every sector of every link type are re-derived and agree with the derived cells. The gain law follows from the map from the signed exchanges a record survives to the dimension of the end each acts at, and from the two per-end tests at the commutants; the three per-end factors are derived at the one link from the pair state and the transport, with the five other permutations of the statuses excluded at named cells. The carrier webs are read as chains at the resolvent of the arrow operator.

\begin{chbridge}
The Born rule assigns to an outcome of a projective measurement the probability $\mathrm{Tr}(\rho P)$, and the normalizing trace of a projector is its rank, the multiplicity of the outcome \cite{Chiribella2011}. The derivation supplies the record web with its link-product measure and the Born count as a dressed measure here, and the Born rule, decoherence as a master equation, and the bosons as quanta of fields are named by correspondence.
\end{chbridge}
\begin{chproof}[finite]
The premises are the axiom with its six definitions (chapter 1-1, the text), the link with its six types and the readability order (chapter 1-2, a conditional theorem), the explicit two-point function with its records (chapter 5-1, a theorem) and the sector-to-ends map at the six types (chapter 1-5, an open result read at its standing). The chain back to the axiom runs through chapter 1-2 and the framework chapter. The procedure. The input is the one (2,3) link with its one hundred seventy-two records, each record a link at the dimensions of the ends it passes, a closed end at dimension one, and the six types with their sectors. For every record the multiplicity is read and compared with the product of the dimensions of its passed ends, one at no end, two at the qubit end, three at the qutrit end and six at both. The four end-set sums are taken and compared with the derived cells of the support-rank result, and at every sector of every type the gains histogram, the Born face and the coherent sum are re-derived at both readings of the sector-to-ends map. The aggregation counts the records and the sectors met. The totals at the checkpoints,
\begin{equation}
m(r) = \prod_{\text{passed ends}} d \ \text{at } 172 \text{ of } 172,\qquad (79,\ 175,\ 205,\ 319) \ \text{the four end-set sums, equal at } 4 \text{ of } 4 ,
\end{equation}
the sectors equal at six of six types at both readings, so the record web carries a link-product measure equal to the Born count at every end set and every type and a half-power measure, every record at the root of its multiplicity, equal to the coherent sum at six of six. So the Born count is a measure and the readout its root, which is the theorem, the two hundred forty-six candidate cells matched at two hundred forty-six. The smallest instance on the page is a record passing the qubit end alone, its multiplicity is two, the dimension of that end, and its readout is the root of two over the register. The end set of the records passing both ends sums to three hundred nineteen, the cap of the (2,3) link.

The counter-tests struck the multiplicity rule to the passed-end count, which moved the face to a parting at a mixed record with the multiplicity two against three. struck the sector-to-ends map at the (2,4) type, which moved the gains histogram and the Born face to a parting there, and struck the text at its gate, which moved every cell. Every control read the same at both engines, thirteen of thirteen, and the multiplicities and the sums were reproduced at a second engine built apart, equal at every record.
\chcert{172}
\end{chproof}
\section{The soft limit as the washed open read}\label{mod:I-3-D534981-w-sftp}
\begin{chwall}[The soft limit as the washed open read]\label{box:I-3-D534981-w-sftp:1}
The wash (decoherence) of the open read at a qutrit-qubit dock (the dock: the photon-electron vertex) contributes the unit to the two-point function at the residue at every body, the same element at every dock, and the residue with the dock read washed is the derived residue of the pole-residue result, so the soft-limit structure, one open read at the dock contributing a universal factor fixed at zero displacement, stands at every such dock of every derived body; the structure parts at the ladder ports of the qubit-qubit type, which the adopted strike strikes beside the docks under one rule and at which no open read exists, and the open result stands at that parting.
\end{chwall}
\chstanding{open result}
\chdeps{Definitions 1.1 to 1.6 and Theorems 18.1 to 18.4 and Open results 18.5 and 18.6}{no theorem beyond it at this release}
The washed open read is built at every admitted web with the derived residues of the pole-residue result re-met, the measurement acts at the listed docks re-met and the washed factor the unit at every docked web; the struck links of every web are compared with the derived cells as unordered end pairs with their measures, agreeing at every one of the seventy-six webs, and the comparison parts at the ladder ports, where the adopted strike strikes beside the docks under one rule and no open read exists.

\begin{chbridge}
In the soft limit of quantum electrodynamics the emission of a photon of vanishing momentum contributes a universal factor to every amplitude, the same at every charged external line, and the infrared divergences of real and virtual soft photons cancel order by order in the Bloch-Nordsieck sum \cite{Low1958,BlochNordsieck1937}. The derivation supplies the two-point function of the one link at zero displacement and the washed read at the dock here, and the photon-electron vertex, the coupling at zero momentum transfer, and the soft limit of an amplitude are named by correspondence.
\end{chbridge}
\begin{chproof}[census]
The premises are the pole-residue result with its derived residues (chapter 5-1, a theorem), the measurement as the wash at one act (chapter 1-6, a theorem) and the admitted webs of the saturation with their docks (chapter 1-5, a theorem). The chain back to the axiom runs through chapter 1-6 and the framework chapter. By exhaustive enumeration of the seventy-six admitted webs, the listed bodies, the mesons and the docked webs, each built with the washed open read at every one of its dock links (the set generated by the strike rule applied at every web, its struck links listed as unordered end pairs with their measures and equal to the derived cells of the pole-residue result at seventy-six of seventy-six). The property tested at every docked web being the soft-limit structure, one open read at the dock contributing a factor fixed at zero displacement and universal over the bodies, with the residue read with the dock washed equal to the derived residue. The counts by kind, the structure equal at twenty-eight of the thirty docked webs, the listed bodies at seven of nine and the mesons at twenty-one of twenty-one, the universality clause at twenty-eight of thirty and the fixed clause at thirty of thirty. The dock links of the (2,3) type at fifty-three of the fifty-six, every one open with the readout equal to the readout of the one link, the Born square the unit and the factor the unit, and the (2,2) ports at three of fifty-six. The first parting element the muon at its port link of the (2,2) type with the far end washed at both orientations, a link the strike strikes beside the docks under one rule and at which no open read exists, which is the open result. The smallest case on the page is the one (2,3) link read at its dock, its common support has dimension two, its readout is one half at each of the two outcomes and its register is six, so the washed open read contributes the factor one,
\begin{equation}
\tfrac12 + \tfrac12 = 1 ,
\end{equation}
and the residue with the dock read washed is the derived residue.

The counter-tests read the docks alone, which moved six cells beginning at the dock count of the muon web. flipped the open end, which moved one hundred ten cells. left the read unwashed, which moved eighty-seven cells. read the dock in place of the open end, which moved no cell as pinned, and struck the text at its gate. Every control read the same at both engines, the struck links of every web equal to the derived cells as unordered end pairs at seventy-six of seventy-six by a second route, and the structure counts were reproduced at a second engine, equal at every cell.
\chcert{76}
\end{chproof}
\section{The read dresses what it does not pass}\label{mod:I-3-D535732-w-rdwp}
\begin{chwall}[The read dresses what it does not pass]\label{box:I-3-D535732-w-rdwp:1}
The one readout rule, that the read dresses what it does not pass, re-derives the unit at the electron, the hub spelling of the lepton ladder at the three leptons together, and the one-dock readout of the cap census at the proton, yet at no one reading of the grid all four; the rule parts at the muon, whose lowest-order count (the first power) at the adopted reading is sixty-nine against the derived ladder cell of nine hundred sixty-three, and the open result stands at that parting with the sub-question of two rules named.
\end{chwall}
\chstanding{open result}
\chdeps{Definitions 1.1 to 1.6 and Theorems 18.1 to 18.4 and Open results 18.5 and 18.6}{no theorem beyond it at this release}
The paths of the measuring read are printed at every body, the register named at the default reading, the pass rule applied at two readings and the first power computed at a grid of readings; the rule re-derives the electron and the lepton ladder at one reading and the proton cell at another, and at no one reading all four, so the elegant solution of the proton and neutron ratios is not found at one clause there; the neutron first power is an output at every reading. The controls agree at both engines except at one listed web spelling a link type reversed, a web entering no body of the result.

\begin{chbridge}
In the Standard Model the charged lepton masses are free parameters, the Yukawa couplings of the lepton doublets to the Higgs field, and the renormalized mass of a lepton is the pole of its dressed propagator, the self-energy summed over the interactions the particle undergoes \cite{Weinberg1967,Dyson1949}. The correspondence is by name.
\end{chbridge}
\begin{chproof}[exact]
The premises are the gain lemma and the Born count (chapter 1-3, a theorem), the link with its cap (chapter 1-2, a theorem) and the saturation at whose sinks the five bodies stand with their measurement webs (chapter 1-5, a theorem). The chain back to the axiom runs through chapter 1-3 and the framework chapter. From the measurement web of each body with the read entered at its joined dock, by the first-arrival read along the entered link (every site at its first-arrival step, the geodesic to the first hub reached the measuring path, the in-legs on it passed and the in-legs off it not passed), we obtain the five measuring paths, the electron, entered along the (2,3) dock at its first site, reaches the hub at the third step with the one in-leg on the geodesic passed. The muon, entered along its dock, reaches two hubs at the third step at once through two in-legs, a tie, both passed. The tau, entered along a (2,2) pendant pair, reaches its hub at the second step through a (2,4) leg with no in-leg on the geodesic and its three in-legs not passed. The proton, entered along the dock at its first site, reaches the core and the hub at the second step through one in-leg passed and leaves two in-legs not passed. The neutron reaches its hub and its core at the second step through one in-leg passed and leaves three in-legs not passed. From the entered link of each body, by the register typing (the cap of the entered (2,3) link at a body whose dock is capped at the sink, the closed count of the link otherwise), we obtain the registers
\begin{equation}
r_{e}=79,\qquad r_{\mu}=79,\qquad r_{\tau}=24,\qquad r_{p}=319,\qquad r_{n}=319 .
\end{equation}
From the passed objects, by the strike of the entered dock (the records removed at the dock, twelve at the electron, ten at the muon, eight at the tau, eighteen at the proton and eighteen at the neutron) and by the dressing of every in-leg not passed by a factor of three, we obtain the five first powers as exact integers,
\begin{equation} \adjustbox{max width=\dimexpr\linewidth-3em\relax}{$\displaystyle p_{e}=(79-12)\cdot 1=67,\quad p_{\mu}=(79-10)\cdot 1=69,\quad p_{\tau}=(24-8)\cdot 27=432,\quad p_{p}=(319-18)\cdot 9=2709,\quad p_{n}=(319-18)\cdot 27=8127 .$} \end{equation}
Against the cells derived elsewhere in this compendium, the readout unit sixty-seven at the electron (chapter 1-2), the ladder cells nine hundred sixty-three at the muon and three thousand nine hundred sixty-nine at the tau, and the candidate cell two thousand eight hundred seventy-one at the proton, the first power is equal at the electron and parts at the muon, the tau and the proton. The first parting body in the order electron, muon, tau, proton is the muon, with the cells sixty-nine against nine hundred sixty-three, which is the open result. The neutron first power is printed as an output. The smallest instance on the page is the electron, seventy-nine records at the entered link less the twelve of the dock strike, no in-leg left unpassed, so the first power is sixty-seven, the unit. Over the grid of sixty readings, four of the register, three of the records and five of the path, no one reading meets the four derived cells at once, the reading at the web measure with the strike removed and the tied legs unpassed meets the three leptons, sixty-seven, nine hundred sixty-three and three thousand nine hundred sixty-nine, and parts at the proton. The proton is met at six readings, every one at the capped register of its entered link with nothing removed,
\begin{equation}
319\cdot 9 = 2871 ,
\end{equation}
so the passed records at the leptons and the passed leg at the proton are objects of two rules, the register a web measure with the strike removed at the leptons and the cap of one link with nothing removed at the proton.

The counter-tests struck the path reading to the orbit reading, the register typing to the cap, the passed records to none and the text of the axiom at its gate, and each moved the cell it was pinned to move. The controls read the same at both engines but one, a listed web whose twelve links spell a link type in the reversed order, a web entering no body of this result. The paths, the registers and the five first powers were reproduced at a second engine, equal at every cell.
\end{chproof}
\begin{figure}[h]\centering
\begin{adjustbox}{max width=\linewidth}\begin{tabular}{llrrrr}\toprule
body & entered link & hub reached at step & in-legs passed & in-legs not passed & register \\ \midrule
electron & a (2,3) dock & 3 & 1 & 0 & 79 \\
muon & a (2,3) dock & 3 & 2 & 0 & 79 \\
tau & a (2,2) pendant pair & 2 & 0 & 3 & 24 \\
proton & a (2,3) dock & 2 & 1 & 2 & 319 \\
neutron & a (2,3) dock & 2 & 1 & 3 & 319 \\ \bottomrule
\end{tabular}\end{adjustbox}
\caption{The measuring path of the read at the five bodies: the link the read enters, the step at which the first hub is reached, the in-legs on the geodesic passed and the in-legs off it dressed, and the register of the entered link.}
\end{figure}
\section{The read stops at the cap}\label{mod:I-3-D535866-w-rsac}
\begin{chwall}[The read stops at the cap]\label{box:I-3-D535866-w-rsac:1}
The stop clause, that the read stops at the cap, re-derives the unit at the electron and nothing else: at the proton the read stopping at the capped dock link dresses the geodesic leg that the candidate cell of the cap census needs undressed, so the local clause reads the hub numerator where the candidate reads the dock numerator; the clause parts at the muon, its lowest-order count six hundred twenty-one against the derived ladder cell of nine hundred sixty-three, and at the proton, eight thousand six hundred thirteen against the candidate cell of two thousand eight hundred seventy-one.
\end{chwall}
\chstanding{open result}
\chdeps{Definitions 1.1 to 1.6 and Theorems 18.1 to 18.4 and Open results 18.5 and 18.6}{no theorem beyond it at this release}
The stop clause is applied at the paths of the measuring read of every body with the register at the default reading and the pass rule at its two readings; the first power is computed at every body and compared with the derived cells; the clause re-derives the electron and parts at the muon and at the proton, so the elegant solution of the proton and neutron ratios is not found at one clause local to the link; the neutron first power is an output at every reading.

\begin{chbridge}
The renormalized mass of a charged particle is read at the pole of its dressed propagator, the self-energy summed over the interactions it undergoes, with the lepton masses free parameters of the Standard Model and the nucleon masses outputs of the strong dynamics \cite{Dyson1949,Weinberg1967}. The derivation supplies the stop of the read at the capped link and the first powers here, and the mass, the spin and the binding of the bodies are named by correspondence.
\end{chbridge}
\begin{chproof}[finite]
The premises are the gain lemma and the Born count (chapter 1-3, a theorem), the link with its cap (chapter 1-2, a theorem), the read that dresses what it does not pass (chapter 1-3, an open result read at its standing) and the saturation at whose sinks the five bodies stand (chapter 1-5, a theorem). The chain back to the axiom runs through chapter 1-3 and the framework chapter. The procedure. The input is the measurement web of each of the five bodies with the read entered at its joined dock, the register of the entered link at the default reading and the pass rule at its two readings. The rule applied is the stop clause, that the read stops at the capped link, so that the legs beyond the stop are dressed and the legs up to it passed. The first power of every body is computed at the sixty readings of the grid, four of the register, three of the records and five of the path, three hundred cells in all, and compared with the derived cells. The aggregation names the bodies equal and the bodies parting in the order electron, muon, tau, proton. The totals at the checkpoints, the first powers at the stop clause and the derived cells beside,
\begin{equation}
\begin{array}{lrrrr}
\text{body} & \text{electron} & \text{muon} & \text{tau} & \text{proton}\\
\text{first power at the stop clause} & 67 & 621 & 432 & 8613\\
\text{derived cell} & 67 & 963 & 3969 & 2871
\end{array}
\end{equation}
equal at the electron and parting at the muon, the tau and the proton, the first parting body the muon, which is the open result, with the neutron printed as an output. At the proton the read stopping at the capped dock link dresses the geodesic leg the candidate cell needs undressed, so the clause reads the hub numerator eight thousand six hundred thirteen where the candidate reads the dock numerator two thousand eight hundred seventy-one. The smallest instance on the page is the electron, the read enters at its dock, stops at the capped link, passes the one leg on its geodesic and dresses nothing, so its first power is the unit sixty-seven, the only compared body the clause re-derives at the typing taken and at the census and all-capped typings alike.

The counter-tests moved the readings of the path, the register and the passed objects one by one, each moving the first parting, the first powers of the muon, the tau and the proton and the two faces of the pass rule, and struck the text at its gate, which moved every cell, seven of seven at their pinned letters. Every control read the same at both engines, seventeen of seventeen, and the first powers were reproduced at a second engine, equal at every cell.
\chcert{300}
\end{chproof}
\section{The Aharonov-Bohm phase as the enclosed record count}\label{mod:I-3-D535921-w-abrc}
\begin{chwall}[The Aharonov-Bohm phase as the enclosed record count]\label{box:I-3-D535921-w-abrc:1}
The relative holonomy of the two paths is trivial at the adopted reading at every two-path web, and the enclosed dark links (the dark link) enter it nowhere: at the adopted transport the linear element every loop defines is the identity on the common support at every one of the forty-five thousand seven hundred thirty-nine two-path webs built at the dark links of the admitted webs, unmoved by the removal of any enclosed dark link; at the matrix reading beside, a path dependence is present by parity alone at the odd loops; every loop element lies inside the matching-phase group of its loop, and the matching freedom moves the loop element, a free datum the text fixes at no corner. The Aharonov-Bohm phase as the enclosed record count is therefore reached by the derivation nowhere at this floor, and the open result stands at the letter trivial.
\end{chwall}
\chstanding{open result}
\chdeps{Definitions 1.1 to 1.6 and Theorems 18.1 to 18.4 and Open results 18.5 and 18.6}{no theorem beyond it at this release}
Every two-path web is built at the dark links of the admitted webs by the admitted moves and by the web's own second path, the enclosed set listed at every web; the holonomy of every loop is computed at the adopted transport, the conjugate on the common support at every link with the identity matching at every corner, and read as a linear element, trivial at every web; the matrix reading is printed beside; every loop element is placed inside the matching-phase group of its loop and the quantum per link type is printed from the phase entry; the matching freedom is sampled and found to move the loop element.

\begin{chbridge}
The Aharonov-Bohm effect: a charged particle carried around a closed loop enclosing a magnetic flux acquires a phase proportional to the enclosed flux although the field vanishes along the path, the phase a nonintegrable factor of the gauge connection and the observable the holonomy around the loop \cite{Aharonov1959,WuYang1975}. The derivation supplies the transport along the links and the loop element of a two-path web here, and the Aharonov-Bohm phase and the enclosed flux are named by correspondence.
\end{chbridge}
\begin{chproof}[census]
The premises are the record phase group at every type (chapter 1-2, a theorem), the moves in one table (chapter 1-4, a theorem) and the admitted webs with their dark links (chapter 1-5, a theorem). The chain back to the axiom runs through chapter 1-2 and the framework chapter. By exhaustive enumeration of the forty-five thousand seven hundred thirty-nine two-path webs (the set generated at the ten thousand nine hundred forty-six dark links of the eight hundred six admitted webs by the duplication and subdivision moves, forty-three thousand eight hundred twenty webs, and by each web's own second path, one thousand nine hundred nineteen webs, the enclosed set of dark links listed at every one). The property tested being the relative holonomy of the two paths at the transport taken, the conjugate on the common support at every link with the identity matching at every corner, read as a linear element on the common support. The counts by kind, the element the identity at forty-five thousand seven hundred thirty-nine of forty-five thousand seven hundred thirty-nine, the webs with a nonempty enclosed set forty-three thousand eight hundred eight and the webs with the empty set one thousand nine hundred thirty-one, every one trivial, the removal of an enclosed dark link moving the element at zero of forty-three thousand eight hundred forty-four removals, and at the matrix reading beside a path dependence present by parity alone at five hundred forty-five webs, the odd loops, every one with the empty set, and absent at forty-five thousand one hundred ninety-four. Every loop element inside the matching-phase group of its loop at every web, the full turn of that group at the orders four, twelve, eight, twelve, twenty-four and eight at the six link types, and the matching freedom moving the loop element at nineteen of nineteen sampled kinds, a free datum the text fixes at no corner. The first parting element the first two-path web with a nonempty enclosed set, a duplicated and subdivided loop of three (3,3) links at a docked electron web enclosing one dark link, whose element is the identity, so the phase as the enclosed count is reached nowhere, which is the open result. The smallest case on the page is the one (2,3) link, its twelve two-path webs are every one trivial, the path dependence present at none and the enclosed dark count zero,
\begin{equation}
h(\gamma) = \mathbb{1}\ \text{at } 12 \text{ of } 12,\qquad n_{\text{dark, enclosed}} = 0 .
\end{equation}

The counter-tests removed an enclosed dark link, which moved the element at no web where the structure sought would have moved it, the finding itself. exchanged the two paths, which reversed every loop element at seventy-two of seventy-two kinds. took the loops enclosing no dark link, trivial at the transport taken and nontrivial by parity at five hundred forty-five at the matrix reading. struck the conjugate at the transport, which moved the cycle row re-met at five hundred eighty-seven of one thousand one hundred eighty-three, and struck the text at its gate. Every control read the same at both engines, and the holonomies were reproduced at a second engine, equal at every web.
\chcert{45739}
\end{chproof}
\section{The geodesic as the read passage and the far-end tie}\label{mod:I-3-D538343-w-gpfr}
\begin{chwall}[The geodesic as the passage and the tie as the far-end rule]\label{box:I-3-D538343-w-gpfr:1}
Read at the five bodies of the base result, the geodesic as the passage of the read and the tie as the far-end rule at the hub legs part at the first body, the electron: the passage set at the adopted measuring path equals the dressed set of the reading alone at the three leptons and is vacuous at the two nucleons, so no one reading is uniform over the five bodies; the tie the far-end rule fixes at the hub legs is neither at the three leptons and both at the two nucleons, agreeing with the recorded tie at four bodies and parting at the muon; and the dressing re-derived with the passage legs dressed and the far-end tie agrees with the derived cells at one body and parts first at the electron, three against one. The open result stands with the first body named.
\end{chwall}
\chstanding{open result}
\chdeps{Definitions 1.1 to 1.6, Theorems 18.1 to 18.4 and Open results 18.5 and 18.6 and Theorems 185.1 to 185.12}{no theorem beyond it at this release}
The passage set is read at the adopted measuring path with the in-legs (the in-leg) typed shared along it and compared with the dressed set at every body; the tie is read at the hub legs by the far-end rule and compared with the recorded tie; the dressing is re-derived with the passage legs dressed and compared with the derived cells under the polarity the item words, the complement polarity printed beside; the two readings part at the electron.

\begin{chbridge}
A free body follows a geodesic of the metric, the path along which the action is stationary, and a measurement passed along an entangled pair fixes the conditional state of the far end in the measured basis and leaves the far marginal unchanged, the no-signalling condition; in the standard account the geodesic and the far-end rule are a statement about spacetime and a statement about measurement with no common premise \cite{Einstein1916,EPR1935,Schrodinger1935}. The derivation supplies the measuring path and the dressing of the unpassed legs here, and the mass and the spin of the five bodies, and the geodesic as a path of a metric, are named by correspondence.
\end{chbridge}
\begin{chproof}[census]
The premises are the measuring path of the read at the five bodies (chapter 1-3, an open result read at its standing), the far-end rule of the read (chapter 1-1, a theorem) and the dressing of the unpassed legs with its derived cells (chapter 1-3, an open result read at its standing). The chain back to the axiom runs through chapter 1-3 and the framework chapter. By exhaustive enumeration of the five bodies of the base result, the electron, the muon, the tau, the proton and the neutron (each at its measurement web with the read entered at its dock, the set generated by the four readings of the passage and the far-end rule at its hub legs, one hundred twenty-six cells in all). The property tested being whether the geodesic is the passage of the read, the passage set at the measuring path equal to the dressed set of one reading, and whether the tie at the hub legs is the one the far-end rule fixes, with the dressing re-derived from the two equal to the derived cells. The counts by kind, the passage set at the measuring path equal to the dressed set of the reading alone at the three leptons and vacuous at the two nucleons, whose measuring path is empty at their instances of no strike, so a reading uniform over the five bodies is found at zero of four and the first body unmatched is the proton. The tie the far-end rule fixes at the hub legs neither at the three leptons, where no dock stands at the far end of any hub leg, and both at the two nucleons, where every hub-leg dock is both, equal to the recorded tie at four of five and parting at the muon, neither at one against both at nine. The dressing re-derived with the passage legs dressed and the far-end tie
\begin{equation}
(3,\,1,\,27,\,1,\,1)\ \text{against the derived cells}\ (1,\,9,\,27,\,9,\,27)
\end{equation}
equal at one body of five and parting first at the electron, three against one, the first power there two hundred one against sixty-seven, with the complement polarity beside at one, one, one, twenty-seven and eighty-one parting first at the muon. The first parting element the electron, at the re-derived dressing three against the derived cell one, which is the open result. The smallest case on the page is the electron, its measuring path carries one in-leg typed shared, so the passage set equals the dressed set of the reading alone, no dock stands at the far end of its hub leg, so the tie is neither, and the one passage leg dressed gives the factor three where the derived cell is one.

The counter-tests re-keyed the passage sets of the muon and the tau and the dressing of the tau, which moved the first parting and the cells of the electron. fixed the tie at the proton and the neutron, which moved the three tie cells. re-keyed the dressing of the electron, which moved nine cells and the first parting, and struck the text at its gate, which moved every cell. Every control read the same at both engines, and the passage, the tie and the dressing were reproduced at a second engine, equal at every body.
\chcert{5}
\end{chproof}
\section{The strong reading of the rate law reconciled}\label{mod:I-3-D538602-rres-t1}
\begin{chconditional}[The three readout rules as one read at three levels]\label{box:I-3-D538602-rres-t1:1}
Conditional on the twelve results this theorem rests on at their own standing, among them the gain lemma, the rate law, the readout unit, the equilibrium theorem with its one-act wash and its wash at the cap, and the base result that names the three rules.
The three adopted readout rules are one read of the junction at the three levels of the wash: at the cap register of the base result the junction record state is carried by the one-act wash to the state beyond one wash and by the wash at the cap to the fixed point, the three states distinct, and each rule equals the junction read at the level the item pairs it with, the gain lemma at the root level, the rate law at the washed level and the readout unit at the flat level; the parting of the base result is therefore a level parting of one read, and the strong-reading open result is superseded with the adopted level named. The result is conditional on the twelve results it rests on at their own grades.
\end{chconditional}
\chstanding{conditional theorem}
\chdeps{Definitions 1.1 to 1.6, Theorems 18.1 to 18.4 and Open results 18.5 and 18.6 and Theorem 42.1}{no theorem beyond it at this release}
The junction record state at the cap register is carried by the one-act wash and by the wash at the cap, and the read summed over every pair of records is evaluated at the root level, the washed level and the flat level, each value re-met against the gain lemma, the rate law and the readout unit; the three level reads are distinct and each rule is met at its level, with the opposite face unmet at every clause.

\begin{chbridge}
A measurement on a quantum system is described at three levels: the amplitude of the outcome ahead of any decoherence, the probability as its squared modulus by the Born rule once the environment has decohered the superposition, and the fixed point of repeated measurement in one basis, the quantum Zeno limit, where the state is pinned to the measured projector; the three descriptions are one process read at three stages of the loss of coherence \cite{Born1926,Zurek2003,Misra1977}. The derivation supplies the wash as a group average and the Born count as a dressed measure at chapters 1-3 and 1-6, and the depolarizing channel and the Kraus decomposition are named by correspondence.
\end{chbridge}
\begin{chproof}[exact]
The premises are the gain lemma and the Born count (chapter 1-3, a theorem), the rate law (chapter 9-2, a theorem), the readout unit (chapter 1-2, a theorem), the wash at one act (chapter 1-6, a theorem), the two washes of the equilibrium theorem (a theorem outside the chapters of this compendium) and the saturated proton with its junction (chapter 1-5, a theorem). The chain back to the axiom runs through chapter 1-3 and the framework chapter. From the junction of the saturated proton, a (3,3) link read at its cap, by the Born count of the end set at both ends, we obtain the register and the junction record state, four hundred twelve records at their multiplicities,
\begin{equation}
1036 = 184\cdot 1 + 128\cdot 3 + 72\cdot 3 + 28\cdot 9 ,
\end{equation}
every record carried at the root of its multiplicity over the register. From that state, by the read summed over every pair of records at the root level, we obtain the root face of the junction read as an exact element of the ring in its basis,
\begin{equation}
J_{\mathrm{root}} = \frac{67}{259} + \frac{50}{259}\sqrt{3} ,
\end{equation}
which is the quantity of the gain lemma. Its square, the plain face at the same level, is $11989/67081 + (6700/67081)\sqrt{3}$, and the unit face, the plain face over the normalized trace, is $47956/259 + (26800/259)\sqrt{3}$. From the junction record state, by the one-act wash (the pinching in the letter of the records, the four hundred twelve record projectors its elements, idempotent as a map and trace-preserving), we obtain the state beyond one wash, diagonal at the register of the records and not flat, whose read at the plain face is the sum of the squared readouts,
\begin{equation}
J_{\mathrm{washed}} = \sum_{\text{records}} \frac{m}{1036^{2}} = \frac{1036}{1036^{2}} = \frac{1}{1036} ,
\end{equation}
which is the quantity of the rate law. The root face at this level is undefined in the field. From the state beyond one wash, by the wash at the cap (the second read, idempotent with the multiples of the identity its image), we obtain the fixed point, the flat state, whose read at the unit face is one,
\begin{equation}
J_{\mathrm{flat}} = 1 ,
\end{equation}
which is the quantity of the readout unit. The three states are distinct at every one of the three pairs and the three level reads are distinct. each rule equals the junction read at the level it is paired with, the gain lemma at the root level, the rate law at the washed level and the readout unit at the flat level, equal at three of three with no parting, so the three rules are one read at three levels of the wash and the parting between them is a level parting of one read. The smallest instance on the page is the (2,3) dock link read at its cap of three hundred nineteen records with one hundred seventy-two records, its root read is $79/319 + (48/319)\sqrt{2} + (42/319)\sqrt{3} + (3/319)\sqrt{6}$, its washed read $1/319$ and its flat read one, the pairing of the rules to the levels equal at three of three there too.

The counter-tests struck the one-act wash and the wash at the cap to the identity map, which carried the three levels to one state and the rate law and the readout unit to a parting at eleven cells. permuted the pairing of the rules to the levels cyclically, which parted every rule. put the dock link in place of the junction, which moved the register and the root and washed reads and left the pairing equal, and struck the text of the axiom at its gate. Every control read the same at both engines, the two engines built apart with no shared line, and the three level reads were reproduced at a second engine, equal at every cell.
\end{chproof}
\begin{figure}[h]\centering
\begin{adjustbox}{max width=\linewidth}\begin{tabular}{lrrr}\toprule
records of the junction & gain & count & weight in the register \\ \midrule
at gain one & 1 & 184 & 184 \\
at gain three & 3 & 128 & 384 \\
at gain three & 3 & 72 & 216 \\
at gain nine & 9 & 28 & 252 \\ \midrule
the register & & 412 & 1036 \\ \bottomrule
\end{tabular}\end{adjustbox}
\caption{The junction of the saturated proton, a (3,3) link read at its cap: its four hundred twelve records at their gains, the register one thousand thirty-six their weighted count.}
\end{figure}
\section{The record algebra at a link}\label{mod:I-3-D539072-w-calr}
For a finite group acting on a finite-dimensional space, the commutant of the action is the algebra of operators commuting with every group element, and the twirl (the average of the conjugation action over the group) is a projector onto that commutant; when the action is a direct sum of inequivalent irreducible representations the commutant is spanned by the projectors onto the irreducible blocks and is commutative \cite{Schur1905,Weyl1931,Werner1989}. The records, the reads and the Born count are introduced at section~\ref{mod:I-3-records-reads-and-the-born-count-the-gain-lemma} and the link with its six genres at section~\ref{mod:I-2-link-and-its-six-genres}.
The end algebra that the known clause lands at a link, the commutant of the matching letters on the identification space of the link, is spanned by the two end projectors, equals the range of the twirl and is commutative at every one of the six genres of the standard links and at the alpha star; the clause that two reads commuting only up to a sign form a further record therefore closes inside it at no record at the strict sign.
At each of the six genres, and at every one of the 1691 standard links, the commutant on the identification space has dimension two and is the span of the two end projectors, the twirl acting as a channel on the matrix units having the same range, so the sign clause generates 0 records at the strict sign at every genre. At the alpha star the closure inside the end algebra generates the 169 degree-two records at none of eight readings, 0 at the principal reading, where the first ungenerated record is a centre-against-far pair, and 6 at the wide sign, none of them among the 169; inside the traceless end algebra the base result is re-derived at 26, 6 and 20 records at the strict sign, reproducing its tables. At the 45739 dark-link entries of the fan table, 10947 dark links over 790 webs, the record count 0 equals the odd-fan count of the enclosed set in half turns at the 12329 entries with no odd fan and at 0 of the 10947 dark links elsewhere; the first link to part is a one-link web at which the record count 0 stands against an odd-fan count of 2, the fan sum 1. The readings beside agree with the odd-fan count at 12326 entries and with the enclosed count of the flux clause at 1386, the fan sum in turns reproducing the fan result at 34215 of the 45739.
The result is an open one, an open result: the stronger statement, that the end algebra generates the degree-two records of the alpha star, is not met at either of its conjuncts, and the weaker statement is named with both values, the first slot parting ahead of the first link parting.
What stays open is a reading at which the end algebra carries the degree-two records of the alpha star under the sign clause, found at none; the record rule at the closure (section~\ref{mod:I-5-D538628-w-lrkc}) and the end algebra of this section await a result that reads them together.
\section{The baseline webs at the loop read of the corridor}\label{mod:I-3-D539202-w-nupr}
A charged quantum passing around a closed loop that encloses a flux acquires a phase fixed by the enclosed flux alone, a holonomy of the connection along the loop, and the interference of the two routes around the loop reads that phase as a shift of the fringe pattern; the phase is a property of the loop as a whole and of nothing local on it \cite{Aharonov1959,Berry1984}. The records, the reads and the Born count are introduced at section~\ref{mod:I-3-records-reads-and-the-born-count-the-gain-lemma} and the detector as a web at section~\ref{mod:I-6-read-at-a-register-the-probe-and-the-measurement}.
The three baseline webs of the oscillation result are rebuilt with the detector as a loop enclosing the corridor and the classical chain, the read being the count of enclosed records and the oscillation the holonomy of the loop. At this stage the loops are unlaid at the four two-route webs and closed at one of the two single-route webs, the reads undefined at the unlaid loops and zero at the closed one, so the loop read of the corridor is open or unlaid at every baseline web and zero where it is laid.
At the reading taken here, the far end a record, the chain site a body and the chain length a record, three of the four two-route webs have no detector body at the base and their chain is laid to no body, while the first web of the second key has its chain laid at no site because the source body stands at its valence; the single-route web of the third key closes, its chain six links long passing three dark links with the enclosed count 0, the holonomy the identity inside a group of order 24 and three copies read, and its companion web is absent. The reads stand undefined at five webs and 0 at one against the through read of the base, which is empty at every body; the holonomy differences between the keys are undefined, and the mixing overlap at the third key is reproduced at both of its entries. Every reading beside, over the component and the joined readings at the record and the least chains, gives the same letters, unlaid at four webs, closed at one and none at the last.
The result is an open one, an open result: the loops are unlaid or read zero at the baseline webs, and the result stands at the weaker reading, its first parting at the first key where the read is undefined against a nonzero target.
What stays open is a baseline web whose corridor route carries a dark chord off the loop or a parallel dark link, found at none; the reading of the oscillation of the base as a loop holonomy awaits a result beyond this stage.
\section{The phase at the junction at the read order}\label{mod:I-3-D540846-phjr-t1}
\begin{chtheorem}[The phase at the junction at the read order]\label{box:I-3-D540846-phjr-t1:1}
No phase at a single (3,3) link of the roster is readable at the adopted order at any sink: at the adopted link of every one of the forty-eight webs carrying a (3,3) link, four hundred ninety-four links in all, every one typed absent at order one and present at order two with the closure twenty-eight, the phase of the shared root of order three multiplied into every amplitude of the pair ket moves no recorded order-two cell at any phase, the Born squares one over eighty-one and sixty-four over six thousand five hundred sixty-one, the sector sums, the weight sum, the strong-coupling cell through the Z body, the office one over eighty-one, the second-order read and the closed read of the word, the moved cells zero in all at forty-eight of forty-eight webs and four hundred ninety-four of four hundred ninety-four links, while the order-one coherent sum of the pair ket carries the phase at forty-eight of forty-eight, the read absent at order one by the confinement row; the block control moves as the quark-mixing-block row prints, the Jarlskog of every recorded cycle equal to the imaginary part of its Bargmann invariant at two hundred eighty-eight of two hundred eighty-eight, and at the adopted block the same phase inserted at the entry moves the Jarlskog at every cycle through the entry, four of four at either orientation. The first face is met and the second face, a readable single-link phase at a sink, is unmet, the first failing clause none.
\end{chtheorem}
\chstanding{theorem}
\chdeps{Definitions 1.1 to 1.6, Theorems 18.1 to 18.4 and Open results 18.5 and 18.6, Theorem 147.1, Theorem 166.1 and Open results 166.2 and 166.3, Theorem 171.1 and Open result 171.2, Theorems 181.1 and 181.2 and Open result 181.3 and Theorem 195.1}{no theorem beyond it at this release}
The (3,3) links of the roster are typed at their two orders from the confinement table; the phase of the shared root is multiplied into every amplitude of the pair ket at the adopted link of every web and every recorded order-two cell is re-read at every phase and compared with the zero phase; the block control recomputes the Jarlskog of every cycle of the washed quark block from the recorded Bargmann invariants and inserts the same phase at the entry of the adopted block.

\begin{chbridge}
The phase of a single amplitude is not an observable; only a rephasing-invariant combination along a closed loop of amplitudes can be, and in the quark mixing matrix the one such invariant is the Jarlskog invariant, the imaginary part of a quartet of matrix elements \cite{Jarlskog1985,Kobayashi1973}. The derivation supplies the phase at a single link and at a block of the washed quark block here, and the mixing phase and the violation of charge-parity symmetry it would fix are named by correspondence.
\end{chbridge}
\begin{chproof}[finite]
The premises are the confinement as the readability of the dark genre (chapter 7-1, a theorem), the gluon pair ket through the scale body (chapter 5-7, a theorem), the quark mixing block (chapter 6-7, an open result), the gain lemma (chapter 1-3, a theorem), the conjugate-letter sorting (chapter 7-7, a theorem) and the body-alone office of every body (chapter 8-2, a theorem). The chain back to the axiom runs through chapter 1-3 and the framework chapter. The procedure. The input is the forty-eight webs of the roster carrying a (3,3) link, four hundred ninety-four links, every one equal to the row of the confinement table found by web and label, the first-order readout absent and the second-order read present with the closure twenty-eight, the sectors of the gain lemma one hundred eighty-four, one hundred twenty-eight, seventy-two and twenty-eight at the gains one, three, three and nine; the adopted link of every web is its first (3,3) link in the order of the roster. From the pair ket on the common support of dimension three, eighty-one pair letters with nine nonzero and eight nonzero and traceless at both ends, by the count faces, we obtain the Born squares
\begin{equation}
\tfrac{1}{81} \quad\text{and}\quad \tfrac{64}{6561} ,
\end{equation}
the weight sum one and the sector sums one over nine, zero, zero and eight over nine, equal to the recorded cells of the gluon-pair-ket row at forty-eight of forty-eight recorded links. From the sectors, by the coherent sum two hundred sixty-eight plus two hundred root three over the Born face one thousand thirty-six with the registers eighty-one and sixty-four, we obtain the order-two read
\begin{equation}
G_{2} = \tfrac{67}{82944} + \tfrac{25}{41472}\sqrt{3} ,
\end{equation}
equal to five recorded cells, and the closed read of the gluon word at its measure twenty-seven, $\tfrac{67}{3072} + \tfrac{25}{1536}\sqrt{3}$. The enumeration rule multiplies the phase of the shared root, the phases zero, one and two of the order three, into every amplitude of the pair ket at the adopted link of every web, re-reads every recorded order-two cell at every phase and compares it with the zero phase; the test at each cell is equality; the aggregation counts the moved cells. The moved cells are zero in all at forty-eight of forty-eight webs and at four hundred ninety-four of four hundred ninety-four links, the strong-coupling cell $\tfrac{1633031687}{25\ 076\ 532\ 510\ 720\ 000} + \tfrac{343}{174142586880000}\sqrt{2}$ standing at every phase, while the order-one coherent sum, the histogram of root exponents over the amplitudes forty-five, eighteen and eighteen at the zero phase, turns at every phase with its squared modulus seven hundred twenty-nine standing. The block control, thirty-two blocks at two hundred eighty-eight cycles, recomputes the Jarlskog of every cycle as the imaginary part of the recorded Bargmann invariant in the ring of the twenty-fourth root, equal to the recorded cell at two hundred eighty-eight of two hundred eighty-eight with its square one over one hundred eight at the mixing pairs of the group, and inserts the phase one at the entry of the adopted block, which moves the Jarlskog at the four cycles through the entry at either orientation, the first cycle from zero, zero, one over nine, zero to zero, one over eighteen, zero, one over eighteen, the five cycles beside standing; so a phase is readable at a block alone and at no single link, which is the theorem. The smallest instance on the page is the pair ket at one link at the phase one: every exponent of its amplitudes stepped by one and every Born square unmoved.

The counter-tests struck the order to one, which compared the coherent sum in place of the Born squares so that the phase read at every web where the read is absent at order one. A second struck the square to the unconjugated square, which carried the phase at the real part of the amplitude squared. A third struck the block phase to a whole-row phase, which left the Jarlskog standing at nine of nine cycles under the rephasing, and a fourth struck the text at its gate, which left the gate unmet. Every control read the same at both engines, fourteen of fourteen, and the cells were reproduced at a second engine, equal at every link.
\chcert{494}
\end{chproof}
\par\medskip\noindent{\itshape Having counted what a read leaves, the thread turns to the moves that carry one web to another.}\label{narr:close:1-3}\par
\section{The read at a register: the probe and the measurement}\label{mod:I-6-read-at-a-register-the-probe-and-the-measurement}
\chdesig{1-6}{}{I-6-read-at-a-register-the-probe-and-the-measurement}{I-6-D536293-w-rbac}
\label{narr:open:1-6}This chapter is where measurement enters. Theorem~28.1 states measurement as the wash at one act, Theorem~28.2 builds the star-pair and one-link apparatus webs, and Theorem~28.5 reads the classical limit as the zero-density register. The open results beside them carry the register to its edges: the mass unit spelling, Open result~28.3, the dimension census of the ledger, Open result~28.4, the large-register limit, Open result~29.1, decoupling as the register count below the cap, Open result~30.1, the probe electron at every measurement web, Open result~31.1, and the register as the reach bounded at the cap, Open result~32.1.
\begin{chtheorem}[Measurement as the wash at one act]\label{box:I-6-read-at-a-register-the-probe-and-the-measurement:1}
This chapter establishes that measurement is the wash (decoherence) at one act, and reads the probe at the register. At every one of the two hundred twenty-eight one-act reads of the network list the readout at the open end, computed as exact arithmetic in the cyclotomic field, equals the reduced state of the pair at the far end at every direction, with the Born square the unit at every read, and is unmoved under every change confined to a washed end; so measurement is the wash at one act and the washed ends are the apparatus.
\end{chtheorem}
\chstanding{theorem}
\chdeps{Definitions 1.1 to 1.6, Theorems 2.1 to 2.3, Open results 2.4 and 2.5, Theorem 2.6, Open results 2.7 and 2.8 and Theorems 2.9 and 2.10 and Theorems 18.1 to 18.4 and Open results 18.5 and 18.6}{the theorems of Chapters 1-6, 3-4 to 3-5, 3-9, 5-3, 5-5, 5-8, 6-2, 6-4, 6-7, 6-10, 7-1 to 7-2, 7-4, 7-6 to 7-8, 8-2, 9-1, 9-3, 10-1, 11-1, 12-2 to 12-3, 12-6, 13-2, 14-1 to 14-2, 14-5 to 14-6, 15-1 to 15-3, 15-6}
\begin{chtheorem}[The star-pair and one-link apparatus webs built]\label{box:I-6-read-at-a-register-the-probe-and-the-measurement:2}
The star-pair and one-link apparatus webs are built at every admitted reading for every measurement of the item.
\end{chtheorem}
\chstanding{theorem}
\chdeps{Definitions 1.1 to 1.6, Theorems 2.1 to 2.3, Open results 2.4 and 2.5, Theorem 2.6, Open results 2.7 and 2.8 and Theorems 2.9 and 2.10 and Theorems 18.1 to 18.4 and Open results 18.5 and 18.6}{the theorems of Chapters 1-6, 3-4 to 3-5, 3-9, 5-3, 5-5, 5-8, 6-2, 6-4, 6-7, 6-10, 7-1 to 7-2, 7-4, 7-6 to 7-8, 8-2, 9-1, 9-3, 10-1, 11-1, 12-2 to 12-3, 12-6, 13-2, 14-1 to 14-2, 14-5 to 14-6, 15-1 to 15-3, 15-6}
\begin{chwall}[The mass unit spelling as a theorem]\label{box:I-6-read-at-a-register-the-probe-and-the-measurement:3}
The mass-unit spelling is a theorem of the derived results at five of the six steps of its derivation and carries its sixth step, the electron entry, at an open result; and at the probe hub leg every ququart hub leg of the body and of the probe reads outward, the count relative to the probe parting first at the muon.
\end{chwall}
\chstanding{open result}
\chdeps{Definitions 1.1 to 1.6, Theorems 2.1 to 2.3, Open results 2.4 and 2.5, Theorem 2.6, Open results 2.7 and 2.8 and Theorems 2.9 and 2.10 and Theorems 18.1 to 18.4 and Open results 18.5 and 18.6}{the theorems of Chapters 1-6, 3-4 to 3-5, 3-9, 5-3, 5-5, 5-8, 6-2, 6-4, 6-7, 6-10, 7-1 to 7-2, 7-4, 7-6 to 7-8, 8-2, 9-1, 9-3, 10-1, 11-1, 12-2 to 12-3, 12-6, 13-2, 14-1 to 14-2, 14-5 to 14-6, 15-1 to 15-3, 15-6}
\begin{chwall}[The dimension census of the ledger]\label{box:I-6-read-at-a-register-the-probe-and-the-measurement:4}
The census of thirty-four derived mass and coupling cells, each restated as a ledger entry with its register or unit beside, fits at twenty-five and at none at nine, the first the proton entry at its hub reading; and the fine-structure register is the product of the site dimensions of no derived configuration, arising instead by the placement (the placement) law as a product of the cut of the pendant and the wholes of the two reads, arithmetic alone.
\end{chwall}
\chstanding{open result}
\chdeps{Definitions 1.1 to 1.6, Theorems 2.1 to 2.3, Open results 2.4 and 2.5, Theorem 2.6, Open results 2.7 and 2.8 and Theorems 2.9 and 2.10 and Theorems 18.1 to 18.4 and Open results 18.5 and 18.6}{the theorems of Chapters 1-6, 3-4 to 3-5, 3-9, 5-3, 5-5, 5-8, 6-2, 6-4, 6-7, 6-10, 7-1 to 7-2, 7-4, 7-6 to 7-8, 8-2, 9-1, 9-3, 10-1, 11-1, 12-2 to 12-3, 12-6, 13-2, 14-1 to 14-2, 14-5 to 14-6, 15-1 to 15-3, 15-6}
\begin{chtheorem}[The classical limit as the zero-density register]\label{box:I-6-read-at-a-register-the-probe-and-the-measurement:5}
The classical limit is the zero-density register: at every admitted web the sign-record density counts the pairs of reads commuting only up to a sign against the register size, the webs at zero density are listed at every reading, the gluon and the Higgs bodies among them, and the semiclassical expansion is the expansion in that density, the first power of every sink re-met as the sum of its terms grouped by the sign-pair count.
\end{chtheorem}
\chstanding{theorem}
\chdeps{Definitions 1.1 to 1.6, Theorems 2.1 to 2.3, Open results 2.4 and 2.5, Theorem 2.6, Open results 2.7 and 2.8 and Theorems 2.9 and 2.10 and Theorems 18.1 to 18.4 and Open results 18.5 and 18.6}{the theorems of Chapters 1-6, 3-4 to 3-5, 3-9, 5-3, 5-5, 5-8, 6-2, 6-4, 6-7, 6-10, 7-1 to 7-2, 7-4, 7-6 to 7-8, 8-2, 9-1, 9-3, 10-1, 11-1, 12-2 to 12-3, 12-6, 13-2, 14-1 to 14-2, 14-5 to 14-6, 15-1 to 15-3, 15-6}
The chain is derived at bytes, the one-act clause at every link type and the definition of open at every entry of the bases table, the readout at the open end of every one-act read is computed exactly and compared with the reduced state of the pair at every direction, and the rate cells of the parent results are re-met. The apparatus webs are built from the pairing move at the pairable set (the pairable set) and the attachment read. The sign rule of the depth-monotonicity result is applied at the admitted bases at every site and at the Weyl letters beside, the density against the register at every web, the zero-density set and the expansion in the density at every sink.

\begin{chbridge}
A measurement of a system in a basis is a projective measurement with the Born rule for its outcomes, a measurement with an apparatus is modelled by coupling the system to a second system and reading the second, and a generalized measurement is a positive operator-valued measure whose elements sum to the identity \cite{Chiribella2011}. The derivation supplies the readout at the open end of a one-act read and the wash of the other ends here, and measurement, collapse, and the Born rule are named by correspondence, the probe and the mass-unit spelling standing at their open results.
\end{chbridge}
\begin{chproof}[census and exact]
The premises are the axiom with its definition of open and the record clause of its text (chapter 1-1, the text), the one-act clause as a theorem of the derived act (chapter 1-1, a theorem), the acts of the listed webs and their bodies (chapter 3-7, a theorem) and the state space from the act outputs (chapter 2-1, a theorem). The chain back to the axiom runs through chapter 1-1 and the framework chapter. The exact part. From a one-act read entered at the near end of a link, by the definition of open at the bytes, the far end of the entered link is open when the near end is the larger and washed otherwise, the near end never open and every other end washed. From the pair ket on the common support read at the open far end, by the reduced state of the pair there, we obtain the readout at every direction as exact arithmetic in the cyclotomic field,
\begin{equation}
\langle P\rangle_{\text{open end}} = \mathrm{tr}\big(P\,\rho_{\text{far}}\big),\qquad \rho_{\text{far}} = \mathrm{tr}_{\text{near}}\,\lvert\psi\rangle\langle\psi\rvert ,
\end{equation}
the washed square over the Born square equal to that reduced state with the Born square the unit. The census part. By exhaustive enumeration of the two hundred twenty-eight one-act reads of the listed webs (the set generated by the acts of the list, seventy-eight with an open far end and one hundred fifty with the far end washed) at their nine thousand eight hundred eighteen directions. The property tested being whether the Born face equals the readout at the open end at every direction and whether the readout is unmoved under every change confined to a washed end, the washed basis re-matched, the washed end re-labeled, the washed end struck and the other links struck. The counts by kind, the Born face the open-end readout at nine thousand eight hundred eighteen of nine thousand eight hundred eighteen directions, the washed ends discarded without moving it at two hundred twenty-eight of two hundred twenty-eight acts, the one-act clause re-met at six of six types, the record clause at eight of eight fragments of the text and the definition of open at forty-eight of forty-eight rows of the bases table, the first parting act none. So measurement is the wash at one act and the washed ends are the apparatus, which is the theorem, the join-kept reading of the open end printed beside at twenty-eight open ends against seventy-eight. The smallest instance on the page is the one (2,3) link read from its qutrit end, the far qubit end is open, the reduced state of the pair ket there is flat,
\begin{equation}
\rho_{\text{far}} = \tfrac12\,\mathbb{1}_{2},\qquad \langle P\rangle = \tfrac12 \ \text{at each outcome},
\end{equation}
the Born square the unit, and the readout is unmoved by any re-matching at the washed qutrit end.

The counter-tests of the derivation stood at their pinned letters and every control held, the gate met at every file, reading and fragment. The readouts at every direction were reproduced at a second engine, equal at every act.
\chcert{9818}
\end{chproof}
\section{The classical limit as the large-register limit}\label{mod:I-6-D534793-w-lasc}
\begin{chwall}[The classical limit as the large-register limit]\label{box:I-6-D534793-w-lasc:1}
The large-register limit meets its classical structure at four of the seven limit entries and parts first at the curvature: the two-point function at the register equals the wave form at every link type, the depth metric equals the additivity at both heads, the sum over records equals the additivity at every link type, every cap and every seeded sink, and the arrival meets its structure, while the Laplacian trace over the ladder of the curvature entry is linear in the rung at every reading where the inverse-square form asks the product of the trace with the square of the rung to be constant, and the Gibbs state parts at the flat-weight clause alone. The open result stands with the first parting named.
\end{chwall}
\chstanding{open result}
\chdeps{Definitions 1.1 to 1.6 and Theorems 28.1 and 28.2, Open results 28.3 and 28.4 and Theorem 28.5}{no theorem beyond it at this release}
The large-register limit is computed at the seven limit entries of the item, the two-point function, the depth metric, the curvature, the sum over records, the Gibbs state, the arrival and the conserved functionals, the classical structure of each is pinned beside it, and the two are compared entry by entry at the adopted readings; the first parting is the curvature, named with both cells.

\begin{chbridge}
The correspondence principle of Bohr: quantum predictions pass to the classical ones in the limit of large quantum numbers, and by the theorem of Ehrenfest the expectation values obey the classical equations of motion where the potential varies slowly over the width of the wave packet \cite{Bohr1920,Ehrenfest1927}. The derivation supplies the register tower, the Laplacian trace over the ladder and the sum over records here, and the curvature, the inverse-square law, and the Gibbs state are named by correspondence, the thermal structures standing at chapters 11-1 to 11-3.
\end{chbridge}
\begin{chproof}[census]
The premises are the gain lemma and the Born count (chapter 1-3, a theorem), the link with its cap (chapter 1-2, a theorem), the saturation with its seeded sinks (chapter 1-5, a theorem) and the conserved functionals of the full move set (chapter 4-6, a theorem). The chain back to the axiom runs through chapter 1-3 and the framework chapter. By exhaustive enumeration of the seven limit entries at their objects (the set generated by the item, the two-point function at the six link types, the depth tower at its two heads, the curvature ladder at its rungs, the sum over records at the six types, the two hundred twenty-eight link-end caps and the two thousand two hundred seventy-seven seeded sinks, the Gibbs state at the caps, the arrival at seven hundred thirty-four tested webs and the conserved functionals of the move matrix). The property tested being whether the large-register limit of each entry equals the classical structure pinned beside it. The counts by kind, equal at four of the seven, the two-point function at the register equal to the wave form at six of six types, the depth metric equal to the additivity at two heads with the register tower geometric at the ratio thirty-six at every step, the sum over records equal to the additivity at six of six types, two hundred twenty-eight of two hundred twenty-eight caps and two thousand two hundred seventy-seven of two thousand two hundred seventy-seven seeded sinks, and the arrival equal to the straight line at seven hundred thirty-three of seven hundred thirty-three increments. parting at three, the curvature, the Gibbs state at the flat-weight clause alone, at one hundred fifty of two hundred twenty-eight caps, and the conserved functionals at the strict reading, where the kernel is trivial. The first parting element in the frozen order the curvature, the Laplacian trace over the ladder is linear in the rung at every one of its three readings, the second differences zero at every rung and the increment per rung the one-link Laplacian trace at eighteen of eighteen connection steps, while the inverse-square form asks the product of the trace with the square of the rung to be the same at every rung,
\begin{equation}
n^{2}\,\mathrm{tr}\,\Delta_{n} = 32,\ 136,\ 324,\ 608,\ 1000,\ 1512 \quad\text{at } n = 1,\dots,6 ,
\end{equation}
the same at zero of five comparisons, which is the open result. The smallest case on the page is the first two rungs, the trace rises from the first rung to the second by one link trace, slope two, so the product with the square of the rung rises from thirty-two to one hundred thirty-six and the inverse-square form is met at no pair.

The counter-tests struck the pairing of the entries to their structures, which moved the comparison and five entries. struck the wave form, which moved the two-point entry. struck the ladder reading, which moved the curvature entry and the limit cells, and struck the text at its gate, which moved every cell. Every control read the same at both engines, and the seven limits were reproduced at a second engine, equal at every cell.
\chcert{2277}
\end{chproof}
\section{Decoupling as the register count below the cap}\label{mod:I-6-D534954-w-dcpl}
\begin{chwall}[Decoupling as the register count below the cap]\label{box:I-6-D534954-w-dcpl:1}
At every body of the sub-web families and at every depth of the one-link-per-tick traversal from the docks, the readout by the register count below the depth on the full web and the readout of the sub-web filled by the depth as a web of its own agree at seventy-nine of the ninety-two body-depth pairs at the hub typing by degree, parting first at the listed two-dock body at depth three, where the full-web route carries the far-end dressing of an in-leg from the tick at which the in-leg is reached and the sub-web route from the tick at which its hub stands at degree three; at the hub typings by dimension and by list line the two routes agree everywhere. The decoupling premise of the alpha correspondence, re-posed at the derivation, parts at one reading only, and the open result stands at the letter coupled.
\end{chwall}
\chstanding{open result}
\chdeps{Definitions 1.1 to 1.6 and Theorems 28.1 and 28.2, Open results 28.3 and 28.4 and Theorem 28.5}{no theorem beyond it at this release}
The registers of the admitted bodies are listed, the depths of the traversal from the docks printed, and at every body-depth pair the readout is computed by the two routes and compared at the three hub typings; the derived first powers of the parents agree with the full-web values at nearly every cell, and the cap waves of every saturation are printed beside.

\begin{chbridge}
The decoupling theorem of Appelquist and Carazzone: the effects of a heavy particle on low-energy observables are suppressed by inverse powers of its mass, so that a theory read below a threshold is a theory of the light fields alone with the heavy fields absorbed into its couplings \cite{AppelquistCarazzone1975}. The derivation supplies the register count below a depth and the readout of the sub-web filled by it here, and the decoupling of a heavy mode is named by correspondence.
\end{chbridge}
\begin{chproof}[finite]
The premises are the gain lemma and the Born count (chapter 1-3, a theorem), the read that dresses what it does not pass (chapter 1-3, an open result read at its standing), the measurement webs of the bodies with their docks (chapter 1-6, a theorem) and the saturation with its seeds and targets (chapter 1-5, a theorem). The chain back to the axiom runs through chapter 1-3 and the framework chapter. The procedure. The input is the eighteen bodies of the sub-web families, the nine charged bodies at their web blocks, the five targets and the four seeds, each traversed one link per tick from its docks, the entry the qubit ends of the struck links. At every depth of every body two readouts are computed, by the first route the register count below the depth on the full web with every link above the depth contributing nothing, and by the second route the readout of the sub-web filled by the depth taken as a web of its own, its hubs typed by their degree in the sub-web. The test applied at each of the ninety-two body-depth pairs is whether the two readouts are equal. The aggregation counts the equal pairs and names the first parting. The totals at the checkpoints, equal at seventy-nine of ninety-two pairs at the hub typing by degree, the first parting the two-dock electron body at depth three, the first route reading one hundred seventeen and the second thirty-nine, because the first route carries the far-end dressing of an in-leg from the tick at which the in-leg is reached and the second from the tick at which its hub stands at degree three in the sub-web, two legs beyond, so at that typing a register above the depth moves the readout below it, which is the open result. At the hub typings by dimension and by the list line the two routes are equal at ninety-two of ninety-two pairs, the decoupling a theorem by construction there. The smallest instance on the page is the electron body at its six depths, the links, the measure, the records struck at the entry, the dressing by each route and the two readouts,
\begin{equation}
\begin{array}{lrrrrrr}
\text{depth} & 0 & 1 & 2 & 3 & 4 & 5\\
\text{links below} & 0 & 2 & 5 & 6 & 9 & 10\\
\text{measure} & 0 & 12 & 39 & 51 & 75 & 79\\
\text{dressing, first route} & 1 & 1 & 1 & 3 & 3 & 3\\
\text{dressing, second route} & 1 & 1 & 1 & 1 & 3 & 3\\
\text{readout, first route} & 0 & 0 & 27 & 117 & 189 & 201\\
\text{readout, second route} & 0 & 0 & 27 & 39 & 189 & 201
\end{array}
\end{equation}
with twelve records struck at the entry at every depth beyond the first, equal at five depths and parting at depth three alone, where $(51-12)\cdot 3 = 117$ against $(51-12)\cdot 1 = 39$.

The counter-tests set the hub reading to the list line in place of the degree, which moved the second-route cells and the first parting. set the entry reading to the core sites in place of the docks, which moved the depths, the sub-webs and the deciding cells, and struck the text at its gate, which moved every cell. Every control read the same at both engines, and the readouts were reproduced at a second engine, equal at every pair.
\chcert{92}
\end{chproof}
\section{The probe electron at every measurement web}\label{mod:I-6-D535118-w-emwp}
\begin{chwall}[The probe electron at every measurement web]\label{box:I-6-D535118-w-emwp:1}
The probe electron, a copy of the listed unit body carried where it stands at its dock alone and placed by the attach move at the lowest open dock of a body part elsewhere, is placed at thirty-one of the sixty-six admitted measurement webs and parts at the rest, the first parting a web whose one part has no dock; where placed, the struck set of the probe equals the records the apparatus part carries of its own at all but one web, and the probe hub equals the listed electron hub by isomorphism with its measure less its dock links the unit, so the renormalisation condition of the derivation extends from the ten webs of the unit-as-scheme result to thirty of the sixty-six. The open result stands at the placement.
\end{chwall}
\chstanding{open result}
\chdeps{Definitions 1.1 to 1.6 and Theorems 28.1 and 28.2, Open results 28.3 and 28.4 and Theorem 28.5}{no theorem beyond it at this release}
The probe is placed at every measurement web under the adopted reading, the strike rule is applied at every placed probe and the struck set compared with the records of the apparatus part, and the probe hub is compared with the listed electron hub by isomorphism; the placement parts at the carrier words alone and at the junction-sealed cores, which offer no body dock.

\begin{chbridge}
The on-shell renormalization condition is stated at a probe: the charge of the electron is defined by the vertex at zero momentum transfer and the mass by the pole of the dressed propagator, the same condition at every process into which the electron enters, so that one renormalization fixes the scheme for all \cite{Dyson1949,BjorkenDrell1965}. The derivation supplies the probe body with its docks and its struck set at every measurement web here, and the renormalization condition and the electron as a quantum of a field are named by correspondence.
\end{chbridge}
\begin{chproof}[finite and census]
The premises are the readout unit as the scheme at the one link (chapter 1-2, a theorem), the measurement as the wash at one act with its webs (chapter 1-6, a theorem), the moves in one table (chapter 1-4, a theorem) and the read that dresses what it does not pass (chapter 1-3, an open result read at its standing). The chain back to the axiom runs through chapter 1-6 and the framework chapter. The procedure. The input is the sixty-six admitted measurement webs at their thirty-four entries with two beside, forty-seven of them built here, thirty-one rebuilt at the second-count table and re-met at their site, link, type and pendant cells, eleven meson webs and five star-pair webs consumed at their site tables, eighty-five instances in all, and the unit body, a copy of the listed electron web. For every web the probe is carried where a copy of the unit body stands at its dock alone and placed otherwise by the attach move at the lowest open dock of a body part. The tests applied at every placed probe are whether its struck set under the strike rule, the two dock links of the probe, equals the records the apparatus part carries of its own, and whether the probe hub equals the listed electron hub by isomorphism with its measure less its dock links the unit. The aggregation counts the webs placed, the webs at which the strike and the hub cells hold, and names the first parting. The totals at the checkpoints,
\begin{equation} \adjustbox{max width=\dimexpr\linewidth-3em\relax}{$\displaystyle \begin{gathered} \text{placed } 31 \text{ of } 66\ (\text{carried } 17,\ \text{placed } 14), \\ \text{struck set equal } 30 \text{ of } 66, \qquad \text{scheme met } 30 \text{ of } 66, \\ 79 - 12 = 67 \text{ at every probe} , \end{gathered}$} \end{equation}
the parting at thirty-five webs, nineteen unbuilt here and sixteen with no open dock of a body part, and at one web with a probe, the neutron star read at transfer, where the located row strikes nothing. The first parting element the web of the W word read alone at its office, whose one part has no dock, the placement cell empty with no probe hub and no struck set, which is the open result. So the renormalization condition extends from the ten webs of the unit-as-scheme result to thirty of the sixty-six. The smallest instance on the page is a web at which the unit body stands at its dock alone, the probe is the copy carried there, its two dock links are its struck set, equal to the records the apparatus part carries, and its hub with the legs (2,4), (2,4), (2,4) and (3,4) and the docks (2,3) and (2,3) is the listed electron hub, the measure seventy-nine less the twelve of its dock links the unit sixty-seven.

The counter-tests replaced the attach reading by the lowest pendant site of any dimension, which moved twenty-five letters and the first parting. replaced the strike reading by the host docks struck with the probe docks, which moved twenty-eight. replaced the unit body by the muon web, which moved fifty-eight, and struck the text at its gate, which left the gate unmet. Every control read the same at both engines, ten of ten, and the placements were reproduced at a second engine, equal at every web.
\chcert{85}
\end{chproof}
\section{The register as the reach bounded at the cap}\label{mod:I-6-D536293-w-rbac}
\begin{chwall}[The register as the reach bounded at the cap]\label{box:I-6-D536293-w-rbac:1}
The extent of the measuring read (the reach), derived at the adopted body-alone typing with every link of the three lepton webs uncapped and every link of the two nucleon sinks capped, and read at the bound, where an uncapped link is crossed and counted at its measure and a capped link met is counted at its cap once, equals the register the strike-rule result used at all five bodies and re-derives the five first powers under its one rule; but the chain is a reading at two steps, the typing of which web the measuring read enters and the bound at which its count stops, each still owing a derivation. The open result stands with the two clauses named.
\end{chwall}
\chstanding{open result}
\chdeps{Definitions 1.1 to 1.6 and Theorems 28.1 and 28.2, Open results 28.3 and 28.4 and Theorem 28.5}{no theorem beyond it at this release}
The adopted definitions are transcribed, the reach of the measuring read of every body is derived from the derived cap results at the body-alone typing and read at the bound, and the derived reach is compared with the register of the strike-rule result and with the five first powers under its one rule; the one sentence, that the first power of a body is what its measuring read counts, stands with its strike corollary derived and its register and dressing corollaries each owing one clause.

\begin{chbridge}
The physical mass of a particle is the pole of its dressed propagator, the sum of the self-energy insertions along its line, and what the self-energy counts is fixed by the scheme in which the propagator is read; in the Standard Model the lepton masses are inputs and the nucleon masses outputs of the strong dynamics \cite{Dyson1949,Weinberg1967}. The derivation supplies the reach of the measuring read and the register it counts here, and the mass of the five bodies is named by correspondence.
\end{chbridge}
\begin{chproof}[finite and census]
The premises are the axiom with its definitions of the passage and of a record (chapter 1-1, the text), the cap of a link with the cap reached and left type by type (chapter 1-2, an open result read at its standing), the read that dresses what it does not pass and the read that stops at the cap (chapter 1-3, two open results read at their standing) and the saturation at whose sinks the five bodies stand (chapter 1-5, a theorem). The chain back to the axiom runs through chapter 1-3 and the framework chapter. The procedure. The input is the five bodies at their webs, every link of the three lepton webs uncapped by the type letters of the cap-reachable result and the tracked electron dock link, every link of the two nucleon sinks capped by the cap census, with the measuring read entered at the dock of each body alone. The reach of the read is derived at the bound, an uncapped link crossed and counted at its measure and a capped link met counted at its cap once and crossed never for the count. The test applied at every body is whether the derived reach equals the register the strike-rule result used and whether the five first powers under its one rule re-derive their cells. The aggregation counts the bodies met. The totals at the checkpoints,
\begin{equation} \adjustbox{max width=\dimexpr\linewidth-3em\relax}{$\displaystyle \begin{gathered} r = (79,\ 117,\ 155,\ 319,\ 319)\ \text{at } 5 \text{ of } 5, \\ p = (67,\ 963,\ 3969,\ 2871)\ \text{at } 4 \text{ of } 4\ \text{with the neutron } 8613 \text{ an output} , \end{gathered}$} \end{equation}
so the reach re-derives the register and the first powers at every body. yet the chain is a reading at two steps, the typing, which web the measuring read of a body enters, applied by the unit, hub-assembly and cap census results and stated as one statement by none, and the bound, the count stopping at a capped link met while the passage continues through it, stated by no definition and no derived result. Of the three corollaries of the one sentence, the strike corollary is derived at the pass line with its cells twelve, ten, eight, zero and zero equal, while the register corollary owes the clause which links a read counts and the dressing corollary owes which legs it dresses, the first corollary owing a clause in the order register, dressing, strike being the register corollary, which is the open result. The smallest instance on the page is the electron, its measuring read enters at its dock, crosses its uncapped links and counts them at their measures to the reach seventy-nine, and under the one rule its first power is sixty-seven, the unit.

The counter-tests struck the proton first power and the electron first power in turn, each moving the reach, the strike and the meeting cells of its body. struck the text at its gate, which moved every cell. re-keyed the shared reach of the muon and the proton, which moved the reach cells of every body, and struck the chain reading, which moved the definition cell, the forced chain and the first corollary. Every control read the same at both engines, nineteen of nineteen, and the reach was reproduced at a second engine, equal at every body.
\chcert{5}
\end{chproof}
\par\medskip\noindent{\itshape With the probe and its register in place, the Part closes on the table that names every object it has used.}\label{narr:close:1-6}\par
\section{The instrument tables of this Part}\label{mod:I-L-the-instrument-tables-at-the-repository}
\chdesig{1-L}{}{I-L-the-instrument-tables-at-the-repository}{I-L-the-instrument-tables-at-the-repository}
\begin{chtable}[The instrument tables of this Part]\label{box:I-L-the-instrument-tables-at-the-repository:1}
This Part lists 1 instrument chapter, printed among the instruments at the end of this compendium as chapter~\ref{mod:I-7-rosetta-table-in-full}, and the table beneath names each one with its key cells.
\end{chtable}
\chstanding{instrument}
\chdeps{Definitions 1.1 to 1.6}{no theorem beyond it at this release}
The listing collects the instrument chapters of this Part with their key cells. Each instrument stands at the repository as its chapter stood at the previous release of this compendium, and its tables are read there.
\begin{chbridge}
The tables listed here are instruments of the derivation, and the listing names them by their titles.
\end{chbridge}
\begin{chproof}[instrument]
The listing in full. For each instrument chapter, printed among the instruments at the end of this compendium, the table gives its title, the result boxes it carries, the standing of each box with its citation, and the count of objects its certificate sentence names where it prints one.
{\small
\begin{longtable}{p{38mm}p{63mm}p{30mm}p{14mm}}
\hline
chapter & results & standing & objects \\
\hline
\endhead
The Rosetta table in full & The Rosetta table of the Chain Theory nouns; theorem & instrument; theorem & -- \\
\hline
\end{longtable}
}
The rule that produced it. Every chapter of this Part whose first standing beyond a placeholder reads instrument is printed among the instruments at the end of this compendium with its tables, and its key cells are the titles of its boxes, its standing lines and its certificate count, read from its text as it stood at the previous release.
No derivation stands in this chapter; the tables stand at the repository.
\emph{Check.} Every file listed was re-hashed at the repository and read unchanged, and the listing was re-cut from the files by a second route.
Nothing stands open at this listing.
\end{chproof}
\label{chend:I-L-the-instrument-tables-at-the-repository}
